% Generated extraction from the Japanese source/combined PDF.
\chapter*{完全理解！篠澤広の機械学習$\cdot$数理最適化ゼミ -1限目-（文字层抽取）}
\addcontentsline{toc}{chapter}{完全理解！篠澤広の機械学習$\cdot$数理最適化ゼミ -1限目-（文字层抽取）}
\begingroup\small
\noindent\textit{来源：docs/完全理解！筱泽广的机器学习$\cdot$数学优化研讨课\textasciitilde{}第一讲\textasciitilde{}/完全理解！篠澤広の機械学習$\cdot$数理最適化ゼミ -1限目-.pdf；页数：65；可抽取文字页：59。}\par
\endgroup
\medskip
\begingroup
\setlength{\parindent}{0pt}
\setlength{\parskip}{0.45em}

% --- PDF page 1 ---
\noindent\textit{[第 1 页没有可抽取文字层。]}\par

% --- PDF page 2 ---
篠澤広に多腕バンディット問題を教わりたい\par
プロローグ\par
やけに暑い5⽉のある⽇、俺は眉間にしわを寄せながらPCの画⾯を眺めていた。表計算ソ\par
フトにときおり数字を打ち込んではそれを全部消し、腕を組んで椅⼦にもたれかかる。そ\par
ういう午後だった。\par
そんな俺の姿を向かいの席に座ったままじっと⾒つめていた篠澤さんが、とうとう⼝を開\par
いた。\par
「きょうのプロデューサーはいつにもまして楽しそう\par
• • • •\par
。どうしたの？」\par
「実は、SNSに掲出する広告のことで悩んでいて。3種類ある広告にどう予算を配分する\par
べきか決めかねていたんです」\par
「ふむふむ。--プロデューサー、授業をしよう」\par
「授業？」\par
「うん。広告への予算配分を、数学的に決定する⽅法について」\par
篠澤さんは机の下の鞄からノートを取り出すと、⽴ち上がって俺の隣に座った。控えめな\par
⾹⽔の⾹りが、静かに篠澤さんの存在を主張する。\par
「まずは問題の定義から⼊りたい。いま、プロデューサーの⼿元には3種類の広告があ\par
る。⽬的はいちばん効果の⾼い広告に予算を多く割り当てること。でも、どの広告がいち\par
ばん効果的かは、まだ不明。あってる？」\par
篠澤さんが⾸をかしげる。\par
「その通りです。通常、広告の効果は表⽰数に対するクリック数で算出されます。今回は\par
SNS広告なので『いいね』や新規フォローの獲得なども含めた\par
エンゲージメント率を⾼め\par
るのが主な⽬標ですね」\par
「おーけー。じゃあ、授業をはじめる、よ」\par
1\par
「本題に⼊る前に、⽤語を整理しよう。いまプロデューサーが悩んでいる問題、広告の予\par
算配分は、多腕バンディット問題として解釈できる」\par
「多腕バンディット問題--腕がたくさんある盗 賊 の話、というわけではなさそうです\par
バンディット\par
ね」\par
1\par

% --- PDF page 3 ---
「ふふ、プロデューサー、おもしろいね。--ここでいうバンディットはスロットマシン\par
のこと。ただし、このスロットマシンにはレバーが複数本あって、しかもレバーごとに当\par
腕\par
たりが出る確率が違う。わたしたちの⽬標は、その中からもっとも当たりが出る確率が⾼\par
いレバーを⾒つけること」\par
「なるほど。広告の種類がレバーの数、エンゲージメントの発⽣が当たりに相当するんで\par
すね」\par
「そのとおり。じゃあ、ここでひとつ問題。3種類の広告からいちばん効果的なものを⾒\par
つけるとき、プロデューサーならどうする？」\par
「素直に考えるなら『全部の広告をしばらく出してみて、その結果を測定する』とかでし\par
ょうか」\par
「たしかに総当たりは有効な戦略。でも、今回みたいに時間も予算も限られている中では\par
そうもいかない。ここで出てくるのが探 索と\par
活 ⽤という概念」\par
ExplorationExploitation\par
「多腕バンディット問題において、探索はまだ当たりの確率がよくわからないレバーを引\par
いてみること。将来より⼤きな利益を得るための情報収集活動ともいえる。活⽤は現時点\par
でいちばんよい結果が出ているレバーを繰り返し引くこと。これは短期的な利益を最⼤化\par
するための⾏動と考えてもいい」\par
「重要なのが探索と活⽤のバランス。広告を例に考えてみよう。もし広告Aを数回出して\par
よい結果が出たとき、バランスが活⽤側に偏っていると、どうなると思う？」\par
「活⽤は短期的な利益を追う⾏動だから……広告BやCを差し置いて広告Aを出し続けるこ\par
とになりますね」\par
「そう。実際に広告Aがいちばん効果的ならそれでいいけど、もしも広告BやCのほうがよ\par
り効果的だった場合、結果的に損をすることになる。探索と活⽤の最適なバランスを⾒つ\par
けることが、多腕バンディット問題の核⼼」\par
「多腕バンディット問題のことはよく分かりましたが--探索と活⽤のバランスはどうや\par
って計算するんですか？」\par
「いい質問」篠澤さんが少し嬉しそうに⾔う。 「次は実際に計算式をみていこう」\par
2\par
「はじめに、多腕バンディット問題を数式で記述してみる。ちょっと難しいけど、我慢し\par
て、ね」\par
そう⾔って、篠澤さんはノートに数式を書き始めた。ペン先が紙⾯の上で踊り、インクが\par
数式へと変換されていく。\par
「わたしたちが選べる選択肢の集まりは集合 として表現される。たとえば広告A、B、C\par
がある場合、 になる。反対に、 の中の各広告は 、または『 に含ま\par
K\par
K = \{A, B, C\} K k K\par
2\par

% --- PDF page 4 ---
れる 』という意味合いを強調するために と表される」\par
「1回の試⾏は時刻 で表される。広告の場合は が1回⽬の閲覧、 が2回⽬の閲\par
覧……という具合だね」\par
「時刻 で実際に選ばれた選択肢は と書く。 は『1回⽬の広告閲覧では広告Aが表\par
⽰された』という意味」\par
「広告に反応があったかは -- 回⽬の閲覧 で『いいね』や新規フォローが得られた\par
かとして表す。反応があれば 、なければ になる」\par
「各広告の真の報酬確率、もといエンゲージメント率は という値で表現する。たとえば\par
広告Aのエンゲージメント率が10\%なら、 と書かれる。ただ は『真の』確率\par
だから、わたしたちが最初から知ることはできない」\par
「わたしたちの⽬標は、総試⾏回数を としたときの総報酬 を最⼤化すること。\par
は『合計する』という意味の記号で、数学では連続する値の累計値を求めるときによく\par
使われる、よ」\par
「合計する、つまり は 、⽇本語で書けば『1回⽬から 回⽬\par
までの反応をすべて⾜し合わせる』ということですね」\par
「さすがはプロデューサー、理解が速くて助かる。問題の説明が済んだから、解法の話に\par
移ろうか。まずは基本的なアルゴリズム、$\epsilon$-greedy法から解説していくね」\par
3\par
「$\epsilon$-greedy法はその名のとおり、現時点でもっとも成績のいい選択肢を貪欲に選び続け\par
greedy\par
る⼿法。でも、それだけだと最適解を⾒逃しちゃうかもしれないから、確率$\epsilon$でランダム\par
にほかの選択肢を試すことで適度に探索も⾏う。$\epsilon$は0から1の間の⼩さな値--たとえ\par
ば、0.1とか--を取る。この値が⼤きければ探索を重視するようになり、⼩さければ活\par
⽤を重視するようになる」\par
「時刻 において、選択肢 を選んだときに得られる平均報酬 の推定値は と書く。\par
\textasciicircum{}は推定値を表すのによく使われるから、覚えておくといい、よ。 の計算式は、こ\par
う」\par
k k $\in$ K\par
t t = 1 t = 2\par
t a\par
t a  =1 A\par
r\par
t t a\par
t\par
r  =t 1 r  =t 0\par
$\mu$\par
k\par
$\mu$  =A 0.1 $\mu$\par
k\par
T\par
r\par
$\sum$t=1\par
T\par
t\par
$\sum$\par
r  $\sum$t=1\par
T\par
t r  +1 r  +2 $\cdots$ +r\par
T T\par
t k $\mu$\par
k\par
(t)$\mu$\textasciicircum{}k\par
(t)$\mu$\textasciicircum{}k\par
(t) =$\mu$\textasciicircum{}k\par
I(a  = k)$\sum$i=1\par
t-1\par
i\par
r  $\cdot$ I(a  = k)$\sum$i=1\par
t-1\par
i i\par
3\par

% --- PDF page 5 ---
「式の中に出てくる はIndicator functionと呼ばれるもので、カッコの中に与えられた\par
指 ⽰ 関 数\par
条件が真なら1を、偽なら0を返す。プロデューサー、 の意味はわかる？」\par
True False\par
「 回⽬の試⾏で選んだ選択肢が なら1、違えば0、という意味ですよね」\par
「正しい。つまりこの式の分⼦ は『過去の各試⾏ において、選択肢\par
が選ばれていたらそのときの報酬 を⾜し合わせる』という意味になって、結果的に\par
『選択肢 から得られた報酬の合計』が求められる」\par
「分母 は『過去の各試⾏ において、選択肢 が選ばれていたら1を⾜し\par
合わせる』という意味になって、結果的に『選択肢 を選んだ回数の合計』が求められ\par
る。ということは、 の計算式で求められるのは--」\par
篠澤さんが⾔葉を⽌め、微笑みながら俺のほうを⾒る。俺の理解度を確かめたいのか、そ\par
れとも俺を困らせたいのか、どちらなんだろうか。そんなことを考えながら、俺は残りの\par
セリフを引き継ぐように話し始める。\par
「--いままでに選択肢 から得られた報酬を選択肢 を選んだ回数で割ったもの--つ\par
まり、選択肢 を選んで報酬が得られた割合」\par
俺もペン⽴てからボールペンを取り、篠澤さんが書いた数式の下に新たな式を書き加え\par
た。\par
「正解。だから、 は『次に選択肢 を選んだときは、経験上このくらいの確率で報\par
酬が得られるだろう』という予測になる。--ところでプロデューサー、選択肢 を選ん\par
だ回数がもしゼロだったら、どうなると思う？」\par
「分母が0になって、ゼロ除算が発⽣するから……計算できなくなってしまいますね」\par
「そのとおり。だから、 を計算するときは⼯夫が必要になる。たとえば計算を始め\par
る前にすべての選択肢を最低1回は選んでみるとか、まだ⼀度も選んでいない選択肢の\par
をすごく⼤きな値に設定しておくとか」\par
「$\epsilon$-greedy法ではこれらの式と変数を使って、こんな基準で⾏動を選択する」\par
1. 0から1までの範囲で乱数を⽣成する。\par
2. もし⽣成した乱数が$\epsilon$未満なら、すべての選択肢 からランダムな選択肢 を選ぶ。\par
3. もし⽣成した乱数が$\epsilon$以上なら、 が最⼤になる選択肢 を選ぶ。\par
I\par
I(a  =i k)\par
i k\par
r  $\cdot$$\sum$i=1\par
t-1\par
i I(a  =i k) i\par
k r\par
i\par
k\par
I(a  =$\sum$i=1\par
t-1\par
i k) i k\par
k\par
(t)$\mu$\textasciicircum{}k\par
k k\par
k\par
(t) =$\mu$\textasciicircum{}k\par
時刻t - 1までに選択肢kを選んだ回数\par
時刻t - 1までに選択肢kから得られた報酬の合計\par
(t)$\mu$\textasciicircum{}k k\par
k\par
(t)$\mu$\textasciicircum{}k\par
(t)$\mu$\textasciicircum{}k\par
K k\par
(t)$\mu$\textasciicircum{}k k\par
4\par

% --- PDF page 6 ---
「3.の操作を数式で書くと となる。 は『与えられた引\par
数から、その後の式の値を最⼤にするような値を選ぶ』操作のこと。今度はこの⾏動原則\par
を広告に当てはめて考えてみよう」\par
1. 広告A、B、Cがあり、 である。\par
2. また、これまでの実績で広告Aのエンゲージメント率がもっとも⾼い。\par
「このとき、⾏動はどのように決定される？」\par
「 ということは10\%の確率で探索を⾏うから、90\%の確率で広告Aを出し、10\%\par
の確率で広告A、B、Cのいずれかをランダムに出す、となるはずです」\par
「うん、ちゃんと理解できてるね。最後に、$\epsilon$-greedy法のメリットとデメリットについて\par
まとめるよ」\par
$\epsilon$-greedy法のメリット：\par
考え⽅が直感的に理解しやすい。\par
実装が⾮常に簡単。\par
$\epsilon$-greedy法のデメリット：\par
探索時の選択肢の選び⽅が完全にランダムなので、明らかに成績の悪い選択肢も同じ\par
確率で選んでしまう。\par
最適な$\epsilon$の値を選ぶのが難しい 。\par
「さて、ここでプロデューサーに質問。もしも$\epsilon$-greedy法のデメリットを解消できるよう\par
なアプローチを考えるとしたら、プロデューサーはどうする？」\par
「そうですね……探索を⾏うときに完全にランダムに選択肢を選ぶのではなく、試⾏回数\par
の少ない選択肢を積極的に選ぶのがよいのではないでしょうか。成績が同じだとしても、\par
すでに何度も試した選択肢を選ぶより、まだ成績が不確かなものを調べるほうが合理的で\par
す」\par
「ナイスアイデア。次はそのアプローチを定式化したアルゴリズム、UCBについて説明す\par
る、よ」\par
4\par
「UCB、Upper Confidence Boundアルゴリズム\par
は『楽観的な推定』に基づいて選択肢\par
を選ぶアルゴリズム。UCBアルゴリズムでは『試⾏回数が少ない選択肢にはより⼤きな報\par
酬を得られる可能性が残っている』と考えて、それぞれの選択肢の推定される報酬\par
に不確実性のボーナスを加えた評価値を選択基準に使う。このボーナスとして使われる値\par
が信頼区間の上限--Upper Confidence Bound」\par
篠澤さんが"Upper Confidence Bound"の⽂字の上に「信頼区間の上限」というルビを書\par
き⾜す。\par
a  =t arg max    (t)k$\in$K $\mu$\textasciicircum{}k arg max\par
$\epsilon$ = 0.1\par
$\epsilon$ = 0.1\par
1\par
(t)$\mu$\textasciicircum{}k\par
5\par

% --- PDF page 7 ---
「信頼区間は『真の値がこの範囲の中に⾼い確率--たとえば、95\%--で含まれるだろ\par
う』と推定される区間を表す、統計学の⽤語。たとえば、あるコインを1回投げて表が出\par
たからといって『このコインの表が出る確率は100\%』と断⾔するのが危険なのはわかる\par
と思う。実際には50\%のコインの出⽬がたまたま偏っただけかもしれない。信頼区間はそ\par
ういった不確実性を数値で表現する⽅法。試⾏回数が少ない選択肢は不確実性が⼤きいか\par
ら信頼区間が広くなり、反対に試⾏回数が多くて確信度が⾼い選択肢の信頼区間は狭くな\par
る。UCBはこの信頼区間の上限値を楽観的な指標として採⽤する、というわけ」\par
「なるほど。UCBを導⼊することによって、いままでの成績がよい選択肢と、試⾏回数が\par
少なく未知の可能性を持っている選択肢の両⽅に選ばれるチャンスが与えられる。つま\par
り、探索と活⽤のバランスがうまく取れる--ということですね」\par
「そうだね。時刻 における選択肢 の評価値 はこんな⾵に計算される 。\par
の計算⽅法は$\epsilon$-greedy法のときと同じ」\par
「記号の意味も説明する、ね。まず、 は時刻 までに選択肢 が選ばれた回\par
数を指す」\par
「 は時刻 のNatural logarithm。詳細は省くけど、ここでは が⼤きくなるほど の\par
⾃ 然 対 数\par
増加率が緩やかになることで、ボーナスと --平均報酬のバランスをとって、ボー\par
ナスが平均報酬を圧倒してしまうのを防ぐ効果がある」\par
「 は探索の積極性を調整するための定数。この式が初めて提案された論⽂にならって、\par
が採⽤されることが多い」\par
「 には不確実性ボーナスの増加率を緩やかにする役割がある」\par
「不確実性ボーナスの計算式を⾃然⾔語に訳すなら--こう、かな」\par
「ところで、この式で評価値を計算するにはひとつ問題がある。プロデューサー、どこが\par
問題か答えて」\par
「はい。こちらも と同様に、選択肢 の試⾏回数がゼロ回の時に が計算でき\par
なくなってしまう問題を抱えています」\par
t k U CB  (t)k\par
2\par
(t)$\mu$\textasciicircum{}k\par
U CB\par
(t) =k\par
(t) +$\mu$\textasciicircum{}k\par
N  (t - 1)k\par
c ln t\par
N  (t -k 1) t - 1 k\par
ln t t t ln t\par
(t)$\mu$\textasciicircum{}k\par
c\par
c = 2\par
$\cdot$\par
Bonus =\par
時刻t - 1までに選択肢kを選んだ回数\par
探索定数c $\times$ 総試⾏回数tの⾃然対数\par
(t)$\mu$\textasciicircum{}k k\par
N\par
(t-1)k\par
c ln t\par
6\par

% --- PDF page 8 ---
「うん、そのとおり。だからUCBアルゴリズムの場合も、最初にすべての選択肢を試すと\par
かの⽅法で評価値を初期化する必要がある」\par
「多腕バンディット問題をUCBアルゴリズムで解くときの⾏動原理は$\epsilon$-greedy法と⼤差な\par
い。時刻 における をすべての選択肢 について計算して、その結果がもっとも\par
⼤きい選択肢が として選ばれる。⼀連の流れを数式にすると\par
となる」\par
「じゃあ、ここで演習問題に挑戦してみよう」\par
1. 広告A、B、Cがあり、 、 である。\par
2. 時刻 までの広告Aの累計表⽰回数は80回（ ）で、平均エンゲージ\par
メント率は50\%（ ）だった。\par
3. 時刻 までの広告Bの累計表⽰回数は10回（ ）で、平均エンゲージ\par
メント率は10\%（ ）だった。\par
「このときの広告AとBの評価値、 と を計算してみて。計算式\par
をきちんと理解するのが⽬的だから、電卓を使ってもいい、よ」\par
篠澤さんがノートの余⽩を指さした。俺はPCの電卓アプリを⽴ち上げ、式を解き始める。\par
「 だから、これを\par
に当てはめると……」\par
「 なので、 ですね。次は広\par
告Bか……」\par
t U CB\par
(k)t k\par
a\par
t a  =t\par
arg max  U CB  (t)k$\in$K k\par
c = 2 t = 100\par
t - 1 N  (100) =A 80\par
(100) =$\mu$\textasciicircum{}A 0.5\par
t - 1 N  (100) =B 10\par
(100) =$\mu$\textasciicircum{}B 0.1\par
U CB  (100)A U CB  (100)B\par
ln 100 $\approx$ 4.605 U CB (t) =k\par
(t) +$\mu$\textasciicircum{}k\par
N  (t-1)k\par
c ln t\par
U CB  (100)A $\approx$   (100) +  $\mu$\textasciicircum{}A\par
80\par
2 $\times$ 4.605\par
= 0.5 +\par
80\par
2 $\times$ 4.605\par
= 0.5 +\par
80\par
9.21\par
= 0.5 +  0.115125\par
$\approx$ 0.1151250.339 U CB  (100) $\approx$A 0.5 + 0.339 = 0.839\par
7\par

% --- PDF page 9 ---
「 だから、 です」\par
「ぱちぱちぱちぱち。よくできました。ここで注⽬してほしいのは、広告Aより成績が劣\par
っている広告Bが、評価値で広告Aを上回っていること。こんなふうに、UCBアルゴリズ\par
ムは単純な成績だけではなく不確実性も考慮に⼊れて判断を⾏う」\par
「次の解法に移る前に、UCBアルゴリズムのメリットとデメリットについてもおさらいし\par
ておこう」\par
UCBアルゴリズムのメリット：\par
$\epsilon$-greedy法における$\epsilon$と違って、探索の度合いを調整するパラメータ の値に推奨値が\par
ある。\par
最適な選択肢を選ばなかったことによる損失が⼩さいことが理論的に⽰されている。\par
UCBアルゴリズムのデメリット：\par
$\epsilon$-greedy法に⽐べると計算が少し複雑。\par
報酬が0から1の範囲に収まらないケース（広告の収益額など）でうまく機能しない 。\par
「ちなみに『最適な選択肢を選ばなかったことによる損失』には後悔という名前がある。\par
Regret\par
⼀部の最適化問題では、利益を最⼤化するかわりに後悔を最⼩化することを⽬標にするこ\par
ともある、よ」\par
「それじゃあ最後にもうひとつ、トンプソンサンプリングの話をして授業を締めくくろ\par
う」\par
5\par
「トンプソンサンプリング\par
はベイズ統計学の考え⽅に基づいたアルゴリズム。プロデュ\par
Thompson Sampling\par
ーサー、ベイズ統計学はわかる？」\par
「少しだけ、ですが。たしか……実測値を使って次の予測値を更新するサイクルを繰り返\par
す、という話だったはずです」\par
U CB\par
(100)B $\approx$   (100) +\par
$\mu$\textasciicircum{}B\par
10\par
2 $\times$ 4.605\par
= 0.1 +\par
10\par
2 $\times$ 4.605\par
= 0.1 +\par
10\par
9.21\par
= 0.1 +  0.921\par
$\approx$ 0.9210.96 U CB  (100) $\approx$B 0.1 + 0.96 = 1.06\par
c\par
3\par
4\par
8\par

% --- PDF page 10 ---
「それだけわかっていれば⼗分。ほんとうはベイズの定理からちゃんと説明したいけど、\par
きょうのところは省略する、ね」\par
「--トンプソンサンプリングでは、 『各選択肢の真の報酬確率 はこのくらいの確率で\par
この値を取りそうだ』という分布を考える」\par
「そして、選択肢を選んで報酬が得られるたびに、その結果を使って確率分布を更新す\par
る。選択肢を選ぶときは、各選択肢の確率分布から値をひとつランダムに取り出してき\par
サンプリング\par
て、そのサンプル値がいちばん⼤きかった選択肢を採⽤する」\par
篠澤さんはノートのページをめくると、そこに低くなだらかなものや⾼く尖ったものな\par
ど、いくつかの⼭の形を描いた。滑らかな曲線。昔からこうした図形を描き慣れているの\par
だろうか--などと、俺が出会う前の篠澤さんの姿を想像してしまう。\par
「まだ試⾏回数が少なくて真の報酬確率についての確信度が低い選択肢は確率分布の範囲\par
が広いから、⼩さな値から⼤きな値まで、いろいろな値がサンプリングされる可能性があ\par
る。反対に試⾏回数が多くて確信度の⾼い選択肢は確率分布の⼭が狭く尖った形になっ\par
て、平均報酬の予測値\par
に近い値がサンプリングされやすい」\par
「確率分布の範囲が広い選択肢のサンプル値は に対して⼤きく上振れする可能性があ\par
る。だから、UCBアルゴリズムと同様に試⾏回数が少ない選択肢にも実績のある選択肢を\par
差し置いて選ばれるチャンスがあるんですね」\par
「そのとおり。ただ、UCBアルゴリズムのように探索回数の少ない選択肢にボーナスを⾜\par
しているわけではないから、試⾏回数が少なくても成績の悪い選択肢はしだいに無視され\par
るようになる。$\epsilon$-greedy法の$\epsilon$やUCBアルゴリズムの\par
のような値を⽤意しなくても、⾃然\par
に最適解へと収束していくのがトンプソンサンプリングの特徴」\par
「じゃあ、実際の式を⾒てみようか。広告への反応の有無みたいに報酬を0か1に分類でき\par
る場合、その問題は『ベルヌーイバンディット』と呼ばれる。この場合、選択肢\par
の真の\par
報酬確率 の確率分布にはベータ分布を使うのが⼀般的。ベータ分布は の\par
ように、 と のふたつの値によって計算される」\par
：選択肢 で報酬1を観測した回数 + 1\par
：選択肢 で報酬0を観測した回数 + 1\par
「初期値として1を⾜しているのは、最初の確率分布を平坦にするため。まだなにも試し\par
ていない段階では\par
がどんな値を取るかわからないから。プロデューサー、これを逆説的\par
に考えるとどうなる？」\par
「どの値も等しく取りうる、と⾔えます」\par
「ご名答」篠澤さんが少しだけ⽬を細めてほほえむ。 「あらためて、トンプソンサンプリ\par
ングの⼿順を整理してみよう」\par
$\mu$\par
k\par
$\mu$\textasciicircum{}k\par
$\mu$\textasciicircum{}k\par
c\par
k\par
$\mu$\par
k Beta($\alpha$  , $\beta$  )k k\par
$\alpha$\par
k $\beta$\par
k\par
$\alpha$\par
k k\par
$\beta$\par
k k\par
$\mu$\par
k\par
9\par

% --- PDF page 11 ---
1. 時刻 での各選択肢 について、現在のベータ分布 からサンプル値 をラ\par
ンダムに⽣成する。\par
2. がもっとも⼤きい選択肢 を選択する。\par
3. 選択した選択肢 を実⾏し、報酬 を観測する。\par
4. 観測結果に基づいて、 のベータ分布を更新する。\par
もし （報酬を得られた）なら、 を と更新する。\par
もし （報酬を得られなかった）なら、 を と更新する。\par
「トンプソンサンプリングについての説明はこれで終わり。--それじゃあ最後に、ちょ\par
っと難しい問題を出す、ね」\par
篠澤さんは少しいたずらっぽく微笑むと、ノートに新しい問題を書き始めた。\par
宿題：ベータ分布を使った広告性能の確信度分析\par
あなたは2週間にわたって広告配信を⾏った。現在の各広告のベータ分布パラメータは\par
以下の通り：\par
広告A：\par
広告B：\par
広告C：\par
問1：このときの各広告のエンゲージメント率の期待値（平均）と分散を計算せよ。な\par
お、ベータ分布 の平均は 、分散は で求められる。\par
問2：広告AとBについて、 「真のエンゲージメント率が広告A > 広告Bである確率\par
」を求めよ。\par
問3：各広告の現在の性能、確信度、今後の探索価値を総合的に考慮し、プロデューサ\par
ーとして明⽇の広告配信戦略をどう決めるかを、理由とともに答えよ。\par
「明⽇の授業までに解いてきて。プロデューサーなら、きっとできるはず」\par
「ふふ……楽しみにしてる、よ」\par
t k Beta($\alpha$  , $\beta$  )k k $\theta$\par
k\par
$\theta$\par
k a  =t arg max  $\theta$\par
k$\in$K k\par
a\par
t r\par
t\par
a\par
t\par
r  =t 1 $\alpha$\par
a\par
t $\alpha$  +a\par
t 1\par
r  =t 0 $\beta$a\par
t $\beta$  +a\par
t 1\par
$\alpha$\par
=A 23, $\beta$\par
=A 77\par
$\alpha$  =A 8, $\beta$  =A 12\par
$\alpha$  =A 5, $\beta$  =A 45\par
Beta($\alpha$, $\beta$)\par
$\alpha$+$\beta$\par
$\alpha$\par
($\alpha$+$\beta$) ($\alpha$+$\beta$+1)2\par
$\alpha$$\beta$\par
P ($\mu$  >A $\mu$  )B\par
10\par

% --- PDF page 12 ---
参考⽂献\par
1. Auer, Cesa-Bianchi, and Fischer: Finite-time Analysis of the Multiarmed Bandit\par
Problem, 2002\par
脚注\par
1. 「この問題への対処法として、 みたいに、$\epsilon$を試⾏回数とともに減少させる$\epsilon$\par
スケジューリングがよく⽤いられる。こうすれば序盤は⼤胆に探索して、終盤は活⽤\par
寄りに収束させられる」\par
2. 「"Finite-time Analysis of the Multiarmed Bandit Problem"で提案されたUCBの計\par
算式にはUCB1、UCB2、UCB1-NORMAL、UCB1-TUNEDの4種類があるけど、今回\par
はUCB1を使う、よ」\par
3. 「こういうときはUCB1のかわりにUCB-NORMALやUCB-TUNEDを使うのが⼀般的」\par
4. 「確率整合法やPosterior Samplingとも呼ばれるよ」\par
$\epsilon$  =t $\epsilon$  /t0\par
11\par

% --- PDF page 13 ---
バドミントンと天気\par
——ルンゲ‧クッタ法——\par
（うにかじき）\par
「プロデューサー。今から体育館に⾏く」\par
黒板の前で教壇に両⼿をつき、⽩髪の少⼥はそう⾔った。今⽇の⼣⽅は⾃室でゆっくりす\par
る予定だったが、いつの間にか⾃分の担当に連れ去られようとしていた。\par
イメージトレーニング\par
平⽇の昼下がり。前⽇の予報通り、地⾯を叩きつける鋭い⾬が朝から降り続いている。こ\par
の時期のうだるような暑さも相まって、今、歩いている廊下は地獄と化している。\par
⼀刻も早く抜け出そうと早⾜で⼀本道を進めば、がら空きの教室が左⽬に映った。素通り\par
しそうになったその部屋には、窓際近くの席に座る少⼥、篠澤広がいる。両腕を交差させ\par
て机に乗せ、ノートパソコンの画⾯をじっと⾷い⼊るように⾒つめている。その姿は、と\par
ても10代とは思えない。⾶び級をして14歳に⼤学を卒業すれば、あのような⾵格になれる\par
ものだろうか。\par
そのようにしばし考えていると、彼⼥が⼿招きしていることに気が付いた。扉を開け、ち\par
ょうどよい温度と湿度に空調された空間に⼊り、開放感を覚える。\par
「こんにちは、篠澤さん」\par
「こんにちは、プロデューサー」\par
扉を閉め、何気ない挨拶を交わし、篠澤さんに近づく。\par
「何か⽤事がありましたか？」\par
「それはこっちのセリフ。教室の外からずっとわたしを⾒てたでしょ」\par
無意識に絵になる座り姿を⾒続けていたらしい。彼⼥の質問に、とくに⽤事もなく⾒てい\par
た、などと答えるのは何となく⼩恥ずかしく思っていると、⽂字や画像が並ぶ画⾯が⽬に\par
⼊った。\par
「いえ、ずっとパソコンを⾒て、何を考えているんだろうと興味を持っただけです」\par
「これのこと？」\par
「そうです。何かのプログラムですか？」\par
「うん。バドミントンのシャトルコックの軌道を計算してた」\par
「またなんでそんなものを？」\par
「今週末、千奈と佑芽で遊びに出かけるんだけど、そのときにバドミントンをすることに\par
なって」\par
12\par

% --- PDF page 14 ---
「バドミントン、できるんですか？」\par
「少なくとも試合にはならない、はず。だから少しでも練習しようと思って。今⽇は⾬だ\par
から、イメージトレーニングの代わりにシミュレーションを作った」\par
かなり遠回りなやり⽅だ。普通の⼈なら、うまいプレイヤーの動作を⾒ながらイメージト\par
レーニングしそうなものを、物理学から指南を受けようとしている。もっとも、ネットに\par
転がっているコツや指導者が、⽴つのにも苦労した体⼒のなさを考慮に⼊れていないこと\par
を、早々に感じ取ったのかもしれないが。\par
「どれがシャトルコックの軌道なんですか？」\par
「この曲線がその軌道。左から右に向かって打ち出した様⼦。  mで横に伸びている\par
太めの直線がバドミントンのコートで、真ん中くらいにある縦に伸びた直線がネット。こ\par
のグラフはコートを真横から⾒たイメージ」\par
「曲線の最後が急なカーブになるところが、シャトルコックらしい綺麗な軌道になってい\par
ますね」\par
「ちゃんと物理的な要素を⼊れてシミュレーションしているからね。まぁ⼤雑把ではある\par
けど」\par
「シャトルコックには⾊々な⼒がかかるけど、主要なのは地⾯に向かう重⼒と、移動を妨\par
げる空気抵抗の2つ。この⼒を運動⽅程式に⼊れると、右の式になる。\par
と について解い\par
たのが、このグラフ」\par
z = 0\par
x z\par
13\par

% --- PDF page 15 ---
シャトルコックの図と数式が記述された⼩さなノートを指差す。英字と記号の羅列と申し\par
訳程度の⽇本語が並ぶその⾛り書きは、理系らしさを思わせる。\par
「複雑な式ですね。 はどんな式になったんですか？」\par
「具体的な数式はないよ」\par
「はい？」\par
「解こうと思えば解ける、かも。でも今はその式を導出していない」\par
彼⼥は⽭盾した発⾔をあっさりと⾔う。\par
「解いていないなら、さっきのシミュレーションはどうやって作ったんですか？軌道がわ\par
からないのに」\par
「無理やりコンピューターに解かせた」\par
「あぁ、今はやりのAIというやつですか」\par
「うーん......ちょっと、いや、結構違う、と⾔っていいかな。電卓を叩いて計算できる⽅\par
法をコンピューターに⾃動で、⾼速にやらせた、の⽅が正しい、と思う」\par
篠澤さんが説明し続けるほど、シミュレーションの原理が分からなくなる。数式を解いて\par
いないのに、シミュレーションができるのはなぜなのか？パソコンを使うのに、どうして\par
電卓の話になるのだろう？\par
「--プロデューサー。午後に⽤事はある？」\par
俺が腑に落ちていないことを察したのか、教室の窓を眺めながら彼⼥は質問した。\par
「いえ、とくには」\par
「--このシミュレーションの中⾝、理解してみたくない？」\par
これはあれだ。1時間ほど付き合ってくれという宣⾔⽂だ。⽬の前に⼤量の数式が整然と並\par
ぶ黒板が容易に想像できる。やや気が乗らないのはあるが、この後にやることも資料整理\par
くらいで、シミュレーション⾃体に興味があるのも嘘ではない。何より、⾃分の担当から\par
教えをいただくなんていう貴重な経験は滅多にないだろう。すぐに地獄の廊下へ戻るより\par
はよっぽど良い。\par
「そうですね。ぜひ理解してみたいですね」\par
「そう。なら」\par
椅⼦から⽴ち上がり、教壇に向かってパソコンを持ち運び、\par
「シミュレーションで利⽤した数値計算の⼿法の1つ」\par
黒板にある余った⽩チョークをつまんで、俺の⽅に向き直り、\par
「ルンゲ$\cdot$クッタ法の概要を解説する」\par
担当アイドルの講義が始まった。\par
x,z\par
14\par

% --- PDF page 16 ---
⾬\par
「いきなりだけど、題材を⾬粒の落下に変えるよ」\par
「⾬粒？」\par
「そう。シャトルコックの軌道を表すには の2軸が必要。⼀⽅、⾬粒は空から地⾯の⼀\par
直線で落下、つまり の1軸のみでシミュレーションができる。シャトルコックの場合より1\par
軸だけ少ない分、よりシンプル。今から解説するルンゲ$\cdot$クッタ法はどちらにも適⽤でき\par
るから、より簡単な⾬粒の⽅がシミュレーションの中⾝が分かりやすい」\par
「題材を変えてシミュレーションがうまくいかない、とかはないのですか？」\par
「ないとは⾔えない。でも、⾬粒の落下は、シャトルコックの軌道と同じように、重⼒と\par
空気抵抗の2つが主要な⼒。似たような状況だから、シミュレーションも似たようになる、\par
と思ってもらえれば⼤丈夫」\par
「なるほど。わかりました」\par
俺が同意した直後、彼⼥は黒板の端まで歩き、チョークを⾛らせ始めた。\par
「プロデューサーは⾬粒がどれくらいの速さで落ちてくるか、知ってる？」\par
「考えたことはないですね。おおよそ  m/sといったところでしょうか」\par
「直径  mmくらいの弱い⾬ならその程度。直径  mmくらいの強い⾬なら  m/s、時速換\par
算で  km。住宅街を⾛っている⾞の速さくらいで落ちてくる［1］ 」\par
「そう聞くと、思ったより速いスピードで落ちてくるんですね」\par
「⼀⽅、⾬は地上から数千メートルも上の空から降ってくる［2］ 。そう考えると、⾞の速\par
さ程度で落ちてくるのは逆に遅いようにも思える。この理由は知ってる？」\par
「落ちている最中に、空気抵抗を受けるからですよね」\par
「そのとおり。⾬粒は落下中、空気抵抗によってブレーキがかかる。この空気抵抗は落下\par
速度の\par
乗に⽐例するから、加速はかなり早く終わってしまう。実際、  mも落ちれ\par
ば、⾬粒の速度はほとんど⼀定になる」\par
x,z\par
z\par
5\par
1 3 8\par
30\par
2 100\par
15\par

% --- PDF page 17 ---
⼝頭でのやり取りが終わると、チョークを置いて俺に向き直った。\par
「今の話を、運動⽅程式を使って書くとこうなる。⾬粒が落下する向きに 軸正をとる。\par
⾬粒にかかる⼒は、下向きに重⼒、上向きに空気抵抗の2つ。これを運動⽅程式に⼊れる\par
と、右の式になる。 は速度の時間微分、つまり加速度のこと」\par
「 のことですね」\par
「両辺を で割って、少し簡単にするよ。⾬粒は⾃然落下、つまり初速度が  m/sで落下\par
し始めるとしよう」\par
パラパラ漫画\par
「この式を解けば、⾬粒の落下速度 がわかる。これをコンピューターで計算する⽅法を\par
導くのが今からの⽬的。けど、この はコンピューターで直接使えないから、少し形を変\par
える」\par
「どういうことですか？」\par
「コンピューターに合わせて、連続的な値である をバラバラにする。たとえば......」\par
ポケットからスマホを取り出し、俺の顔にカメラのレンズをスッと近づけ、そのまま数秒\par
じっとしていた。\par
「何をしているんですか？」\par
「プロデューサーの動画を撮ってた。思い出に残そうかと」\par
「今すぐ消してください」\par
「冗談。今から使うよ」\par
そういうと、こちらに歩いてきて、先ほど撮った動画を⾒せてきた。\par
「この動画に映るプロデューサーは滑らかに動いているように⾒える。けど、この動画の\par
正体は、1秒間に30枚の写真が⾼速で切り替わっているだけ。原理はパラパラ漫画と⼀緒」\par
「それが今の説明となんの関係が？」\par
「プロデューサーは動画を取り始めてから 秒後に動いていたと思う？」\par
「それは当然動いているでしょうね」\par
「そうだね。でもこの動画にはその 秒後のプロデューサーは映っていない。このパラパ\par
ラ漫画の写真は 秒間隔で映し出されるから、 秒後の写真は存在しない」\par
z\par
v˙\par
ma =F\par
m 0\par
=v˙ - kv +2 g\par
v\par
v\par
v\par
60\par
1\par
60\par
1\par
30\par
1\par
60\par
1\par
16\par

% --- PDF page 18 ---
いまいち話の繋がりが⾒えてこない中、篠澤さんは俺のそばから離れる。\par
「基本的に、現実世界で動いている物体をコンピューターで記録すると、パラパラ漫画の\par
要領でしかデータが保存されない」\par
黒板に向き直り、新しい図を描き始めた。\par
「⾬粒の落下速度も、パラパラ漫画と同じようにコンピューターで計算させなきゃいけな\par
い。具体的には、 秒間隔で⾬粒の速度だけメモしていく感じ。 秒後の速度は計算しな\par
い」\par
「つまり、計算に必要なところだけ記憶させて、それ以外はすべて無視する、というよう\par
なイメージですか？」\par
「だいたいそんな理解で問題ない」\par
「しかし、そんな⾵に無視して、きちんと現実の速度を計算できるんですか？」\par
「うーん、ちょっと考え⽅が違う、かな。記憶させる速度が、ちゃんと現実の速度と近い\par
値になるように、なんとかして計算する、というイメージの⽅が近い」\par
カタッと、乾いた⼩さな⾳が⽴つ。\par
「後の説明のために、コンピューターに計算させる速度の値に名前を付けておく。⾬粒が\par
落下し始めてから\par
秒後ずつの時刻 を\par
とする。時刻が のときの速度を とする。ただし、 は\par
初速度なので と分かっている。だから、コンピューターに計算させる値は、\par
だけ」\par
定積分の近似式\par
「コンピューターに計算させる値が決まったから、実際に計算させる⽅法を考えていくよ」\par
「ようやく本題に⼊れるわけですね」\par
「そう。でもまだ2割くらいしか説明してないけどね」\par
まだまだ彼⼥の貴重なお話は続くようだ。\par
h\par
2\par
h\par
h 0,h, 2h, 3h, ......,ih, (i+ 1)h, ......\par
t  ,t  ,t  ,t  , ......,t  ,t  , ......0 1 2 3 i i +1 t\par
i v\par
i v\par
0\par
0\par
v  ,v  ,v  , ......,v  ,v  , ......1 2 3 i i +1\par
17\par

% --- PDF page 19 ---
「いったん、最初の⽅程式を⾒てみよう。 は の関数だから、 と置いたよ」\par
「プロデューサーなら、この数式から の解を出すとしたら、どうする？」\par
「......⾒当もつきません」\par
「なら、少し話がそれるけど、この式の ならどうする？」\par
「両辺を で積分すればいいんですよね」\par
「そのとおり。具体的に計算式を書き下せばこう。 は積分定数」\par
「運動⽅程式の⽅も同じようにできる。ここから、 と の関係が導ける」\par
＊ ＊ ＊\par
は を で微分したもの。だから、右辺の を で積分すれば、 が求まる。\par
の原始関数の1つを と定義すれば、 は次のように書ける。 は の関数だか\par
ら、原始関数は の関数で問題ない。\par
は とすれば求まる。\par
-kv +2 g t f(t)\par
=v˙ - kv +2 g =f(t)\par
v\par
y\par
=dx\par
dy x\par
x\par
C\par
y = xdx =$\int$\par
+2\par
x2\par
C\par
v\par
i v\par
i+1\par
v˙ v t f(t) t v\par
v = f(t)dt$\int$\par
f(t) F (t) v v t\par
t\par
v =F (t) +C\par
v  ,v\par
i i +1 t =t  ,t\par
i i +1\par
18\par

% --- PDF page 20 ---
この2つの式から を消去すると、 と の関係式が得られる。最後は、 が の\par
原始関数であることを使ったよ。\par
＊ ＊ ＊\par
「......右辺の積分はどうやって計算できるんですか？」\par
「この定積分は正確に計算することが難しい。だから、この定積分に近い値、近似値を計\par
算する⼿法を編み出す」\par
「......これが、"近い値になるように、なんとかして計算する"という意味ですか」\par
「そういうこと。完璧な落下速度の値を計算することは、この時点できっぱり諦める」\par
そう⾔ってから、篠澤さんは再度、俺に質問を投げた。\par
「プロデューサーにもう1つ質問。数学で定積分を使うとき、プロデューサーはどんな問題\par
をよく解いていた？」\par
「ただの計算問題とか......あとは⾯積の計算とかです」\par
「よく覚えてるね。そう。定積分は⾯積の計算に使える。今回の定積分の値は、この部分\par
の⾯積と⼀緒になる」\par
\{v  =F (t  ) +Ci i\par
v  =F (t  ) +Ci+1 i+1\par
C v\par
i+1 v\par
i F (t) f(t)\par
v\par
i+1 =F (t  ) +v  -F (t  )i+1 i i\par
=v\par
+ F (t)\par
i [ ] t\par
i\par
t\par
i+1\par
=v  +  f(t)dti $\int$\par
t\par
i\par
t\par
i+1\par
19\par

% --- PDF page 21 ---
「定積分を近似することは、この⾯積をなるべく正確に計算することと同じ。その計算⽅\par
法を2つ紹介するよ［3］ 」\par
＊ ＊ ＊\par
求める⾯積は、次のAの⾯積から、Bの⾯積を引いたものに等しい。\par
Aの⾯積 は、台形の⾯積公式を使えば簡単に求まる。⾒た⽬が煩わしくなるから、\par
を とする。\par
Aの⾯積は元の定積分の⾯積に近そう。だから、 で定積分を近似することができる。こ\par
の近似は台形則と⾔うよ。これが1つ⽬の計算⽅法。\par
Bの⾯積 を計算できれば、定積分の⾯積をもっと上⼿に計算できそう。ここで、Bを真\par
っすぐ下ろしてみよう。(正確には、 から台形の斜め線の⽅程式を引く。 )\par
S\par
A\par
f(t  ),f(t  )i i +1 f  ,f\par
i i +1\par
S  =A\par
2\par
h(f  +f  )i i +1\par
S\par
A\par
f(t)dt $\simeq$$\int$\par
t\par
i\par
t\par
i+1\par
2\par
h(f  +f  )i i +1\par
S\par
B\par
f(t)\par
20\par

% --- PDF page 22 ---
下ろした後のBを囲む曲線は、 の3点を通る⼆次関数で\par
近似できそう。そこで、この⼆次関数を使ってBの⾯積を近似しよう。曲線は\par
を通るから、定数 を使って⼆次関数の⽅程式は次のようになる。\par
点 を通るように、 を決めるよ。\par
これで⼆次関数が求まった。Bの⾯積 が次のように近似できるね。\par
Bの⾯積をAの⾯積から引いて、定積分の近似値が次のように計算できる。この近似はシン\par
プソン則と⾔うよ。これが2つ⽬の計算⽅法。\par
(t  , 0), (t ,f  -i i c\par
), (t  , 0)2\par
f  +f\par
i i+1\par
i+1\par
(t  , 0), (t , 0)i i +1 $\alpha$\par
y =$\alpha$(t -t  )(t -i t  )i+1\par
(t  ,f  -c c\par
)2\par
f\par
+f\par
i i+1 $\alpha$\par
→\par
f\par
-  =$\alpha$(t\par
-t\par
)(t\par
-t\par
) = -$\alpha$\par
c 2\par
f\par
+f\par
i i +1\par
c i c i +1 4\par
h2\par
$\alpha$ = -  f  -  =\par
h2\par
4 ( c 2\par
f  +f\par
i i +1 ) h2\par
2(f  - 2f  +f  )i c i +1\par
S\par
B\par
S\par
B $\simeq$ -  $\alpha$(t -t  )(t -t )dt$\int$\par
t\par
i\par
t\par
i+1\par
i i +1\par
=$\alpha$\par
6\par
h3\par
=\par
h2\par
2(f  - 2f  +f  )i c i +1\par
6\par
h3\par
=\par
3\par
h(f  - 2f  +f  )i c i +1\par
f(t)dt$\int$\par
t\par
i\par
t\par
i+1\par
=S  -S\par
A B\par
$\simeq$  -\par
2\par
h(f  +f  )i i +1\par
3\par
h(f  - 2f  +f  )i c i +1\par
=\par
6\par
h(f  + 4f  +f  )i c i +1\par
21\par

% --- PDF page 23 ---
＊ ＊ ＊\par
「--以上で台形則とシンプソン則の2つの近似式を紹介したよ。Bの⾯積を引いている\par
分、シンプソン則は台形則より精度が⾼い」\par
近似精度の⾒積もり\par
「篠澤さん。質問いいですか？」\par
「どうしたの？プロデューサー」\par
「シンプソン則が台形則より精度が⾼そうなのは何となくわかりますが、どのくらい精度\par
が上がったんですか？」\par
「いい質問。ちょうど今からその話をするよ」\par
「お願いします」\par
「近似精度を具体的に⽐較するときには、テイラー展開という道具をよく使う。テイラー\par
展開っていうのは、 を の多項式で展開する⽅法のことで、次のような式。\par
は、 を で 回微分した後に、 を代⼊した値」\par
「本当にこんな展開できるんですか？」\par
「できない関数もあるけど、今はこの展開ができると思って⼤丈夫。この展開ができる理\par
由は長いから省略するね」\par
「わかりました」\par
「この数式を使うと、さっきの定積分も擬似的に計算できる」\par
f(t) t f (t  )(n) i\par
f(t) t n t\par
i\par
f(t) =f  +a  (t -t  ) +a  (t -t  ) +a  (t -t  ) +a  (t -t  ) + ......i 1 i 2 i 2 3 i 3 4 i 4\par
a  =\par
n n!\par
f (t )(n)\par
i\par
=\par
=\par
f(t)dt$\int$\par
t\par
i\par
t\par
i+1\par
f  +a (t -t  ) +a\par
(t -t\par
)$\int$\par
t\par
i\par
t\par
i+1\par
i 1 i 2 i\par
2\par
+a  (t -t  ) +a  (t -t  ) + ......dt3 i 3 4 i 4\par
f  h +  h +  h +  h +  h + ......i 2\par
a\par
1 2\par
3\par
a\par
2 3\par
4\par
a\par
3 4\par
5\par
a\par
4 5\par
22\par

% --- PDF page 24 ---
「この式は、近似値ではなく正確な値なんですよね？」\par
「そう。次に台形則とシンプソン則で計算された値を計算する。先に を計算してお\par
こう」\par
「台形則とシンプソン則を計算するよ」\par
「2つの近似式と正確な値を⽐較すると、台形則は の項から、シンプソン則は の項か\par
らずれていることがわかる」\par
「......」\par
「ふふ......プロデューサー。だいぶ苦しそう」\par
今までと⽐べ物にならないくらい長い式に、俺は⽬眩がしそうになる。ずっと押し黙って\par
黒板を凝視する姿を、彼⼥はやや羨ましそうに⾒ていた。\par
「細かい計算は追えていないのですが......2つの近似式の精度が違うのはわかります。Bの\par
⾯積を近似するだけで、\par
、 の項まで⼀致するのはすごいですね」\par
「うん。しかも⼆次関数っていう、すごい単純な曲線で近似しただけなのにね」\par
「でも、これシンプソン則の⽅が精度が悪いように⾒えますが......台形則なら から、シ\par
ンプソン則なら からずれてくるのならば、シンプソン則の⽅がズレが⼤きくなるので\par
は？」\par
「シミュレーションでは、基本的に に設定する。そうすると、 になるから\par
シンプソン則の⽅がズレが⼩さくなる」\par
「 はでたらめに決めてはいけないと？」\par
「そうなるね。ということで、シンプソン則の⽅が台形則よりも だけ精度が⾼い、とい\par
うのが、プロデューサーの質問に対する回答」\par
f\par
,f\par
c i +1\par
f\par
c\par
f\par
i+1\par
=f  +  h +  h +  h +  h + ......i 2\par
a\par
1\par
4\par
a\par
2 2\par
8\par
a\par
3 3\par
16\par
a\par
4 4\par
=f  +a  h +a  h +a h +a  h + ......i 1 2 2 3 3 4 4\par
2\par
h(f  +f  )i i +1\par
6\par
h(f  + 4f  +f  )i c i +1\par
=f  h +  h +  h +  h +  h ......i 2\par
a\par
1 2\par
2\par
a\par
2 3\par
2\par
a\par
3 4\par
2\par
a\par
4 5\par
=f  h +  h +  h +  h +  h + ......i 2\par
a\par
1 2\par
3\par
a\par
2 3\par
4\par
a\par
3 4\par
5a\par
4 5\par
h3 h5\par
h3 h4\par
h3\par
h5\par
h < 1 h <5 h3\par
h\par
h2\par
23\par

% --- PDF page 25 ---
の近似式\par
「シンプソン則をもう1回⾒てみよう」\par
「この右辺、コンピューターが計算できる式だと思う？」\par
「できないと思います。 は ですが、コンピューターは での速度を記憶できないの\par
で」\par
「そう。でも、なんとかしてこのシンプソン則を使えるようにしたい。これを叶えるため\par
に、 を\par
を⽤いて近似していく」\par
「 だけでなく、 もですか？」\par
「うん。計算するときに、 も近似しておくと、コンピューターに計算できる形を導け\par
る」\par
「シミュレーションができるようになると」\par
「そういうこと。近似するときに、先に計算した のテイラー展開を使う」\par
「ここで、 を次のように、 までの項で近似する」\par
「⼤胆に削りますね。これでも良いんですか？」\par
「問題ない。この近似式をシンプソン則にそのまま適⽤すると、次のようになる。この式\par
は、定積分の正確な値と の項まで合っている」\par
f  ,f\par
c i +1\par
f(t)dt $\simeq$$\int$\par
t\par
i\par
t\par
i+1\par
6\par
h(f  + 4f  +f  )i c i +1\par
f\par
c f(t  )c t\par
c\par
f  ,f\par
c i +1 v  ,f\par
i i\par
f\par
c f\par
i+1\par
f\par
i+1\par
f  ,f\par
c i +1\par
f\par
c\par
f\par
i+1\par
=f\par
+\par
h +\par
h +\par
h +\par
h + ......i 2\par
a\par
1\par
4\par
a\par
2 2\par
8\par
a\par
3 3\par
16\par
a\par
4 4\par
=f  +a  h +a  h +a h +a  h + ......i 1 2 2 3 3 4 4\par
f  ,f\par
c i +1 h3\par
f\par
c\par
f\par
i+1\par
$\simeq$f  +  h +  h +  hi 2\par
a\par
1\par
4\par
a\par
2 2\par
8\par
a\par
3 3\par
$\simeq$f\par
+a\par
h +a\par
h +a\par
hi 1 2 2 3 3\par
h4\par
24\par

% --- PDF page 26 ---
「もともとシンプソン則は の項までしか合わないから、この近似は元のシンプソン則の\par
精度を崩さずに近似できていることになる。だから は までの近似で⼤丈夫」\par
「なるほど」\par
「あとは が で表されればいいね。ここは の定義から直接求めるよ。まず\par
は を計算する。途中、 であることを使っているよ。あと、 は と\par
同じ」\par
「これを使えば、 が次のようになる」\par
$\simeq$6\par
h(f  + 4f  +f  )i c i +1 f  h +i\par
h +2\par
a\par
1 2\par
h +3\par
a\par
2 3\par
h4\par
a\par
3 4\par
h4\par
f  ,f\par
c i +1 h3\par
a  ,a  ,a\par
1 2 3 v  ,f\par
i i a\par
n\par
f (t)(n) =v˙ f(t)  f˙ f (t)(1)\par
f (t)(1) =  (-kv +g)dt\par
d 2\par
=\par
-kv +g $\cdot$\par
dv\par
d ( 2 ) dt\par
dv\par
= -2kvf\par
f (t)(2) =\par
(-kv +g)dt2\par
d2\par
2\par
=  (-2kvf )dt\par
d\par
= -2k( f +v )v˙ f˙\par
= -2k(f - 2kv f)2 2\par
f (t)(3) =  (-kv +g)dt3\par
d3\par
2\par
=  \{-2k(f - 2kv f)\}dt\par
d 2 2\par
= -2k\{2f  - 2k(2v f +v  )\}f˙ v˙ 2f˙\par
= -8k (-2vf +kv f)2 2 3\par
a  ,a  ,a\par
1 2 3\par
$\{$\par
$\{$\par
$\{$a\par
1\par
a\par
2\par
a\par
3\par
=  = -2kv  f\par
1!\par
f (t  )(1)\par
i\par
i i\par
=  = -k\{(f\par
) - 2k(v\par
) (f\par
)\}2!\par
f (t  )(2) i\par
i 2 i 2 i\par
=  = -  k \{-2(v  )(f\par
) +k(v\par
) (f\par
)\}3!\par
f (t\par
)(3) i\par
3\par
4 2 i i 2 i 3 i\par
25\par

% --- PDF page 27 ---
「⾒事に と のみで が計算できますね」\par
「そう。この近似式を使えば、 が だけで計算できる。これで、定積分は\par
の2つだけを使ったシンプソン則によって求められる」\par
シミュレーションの主 軸\par
「篠澤さん。近似値の話が続いて、どうシミュレーションに繋がるかが⾒えないのです\par
が......」\par
「ごめん。いったん流れを整理するよ」\par
気づけば、黒板の7割がテイラー展開と近似式で埋まってしまっている。最初に書いた⾬粒\par
の絵なんて、この式が占める領域の半分にも満たないサイズで端に描かれている。\par
「まず、運動⽅程式があって」\par
「これを定積分を使った式に変えた」\par
「最初は、シンプソン則を使って右辺の定積分を近似して」\par
「 を だけを使ってさらに近似した」\par
v\par
i f\par
i a  ,a  ,a\par
1 2 3\par
f  ,f\par
c i +1 v  ,f\par
i i v  ,f\par
i i\par
=v˙ - kv +2 g =f(t)\par
v  =i+1 v +i\par
f(t)dt$\int$\par
t\par
i\par
t\par
i+1\par
v  $\simeq$i+1 v  +i\par
6\par
h(f  + 4f  +f  )i c i +1\par
f  ,f\par
c i +1 v  ,f\par
i i\par
→\par
$\{$\par
$\{$\par
$\{$a\par
1\par
a\par
2\par
a\par
3\par
=  = -2kv  f\par
1!\par
f (t  )(1)\par
i\par
i i\par
=  = -k\{(f  ) - 2k(v\par
) (f\par
)\}2!f (t  )(2) i\par
i 2 i 2 i\par
=  = -  k \{-2(v  )(f\par
) +k(v\par
) (f\par
)\}3!\par
f (t\par
)(3) i\par
3\par
4 2 i i 2 i 3 i\par
\{f\par
c\par
f\par
i+1\par
$\simeq$f\par
+\par
h +\par
h +\par
hi 2\par
a\par
1\par
4\par
a\par
2 2\par
8\par
a\par
3 3\par
$\simeq$f  +a  h +a  h +a  hi 1 2 2 3 3\par
26\par

% --- PDF page 28 ---
「この の近似によって、シンプソン則は だけで表せるようになった。これが\par
今までの話」\par
不必要な式がごっそりと消され、黒板の6割くらいが空⽩になった。\par
「やりたかったこと⾃体は、割とシンプルですね」\par
「計算が複雑なだけで、ただただシミュレーションできるように近似していたというだけ\par
なんだよね」\par
「それで、これをどうやってシミュレーションに落とし込んでいくんですか？」\par
「えっと、もうこの式で⾬粒の落下速度をシミュレーションできる」\par
「どうしてですか？」\par
「まずはじめに⾔っておきたいのは、 が で計算できるということ。\par
だからね。つまり、シンプソン則は の1つだけで計算できてしまう。このこ\par
とを踏まえると、 の近似値は、 だけで近似できる」\par
「ここで とすれば、 （ のみの式）になる。 は初速度で  m/sと最初から与\par
えられているから、 が求まる。次に にすれば、 （ のみの式）になるから、\par
さっき計算した を使って、 がわかる」\par
「ということは、 と順々に⼊れていけば、 と芋づる式に値\par
が分かっていくと」\par
「そう。プロデューサー、飲み込みが早いね。これで のすべての値が計算\par
できる。ちなみに、これが を近似した理由になっている。普通のシンプソン則の式\par
は、 と が からは求まらない。だから両⽅とも近似する必要があった」\par
ルンゲ‧クッタ法\par
「篠澤さんのシミュレーションは、この式でやってたんですね」\par
「いや、わたしのはこれでやってない」\par
「違うんですか？」\par
「うん。というのも、この計算⽅法はやや不便。 の値を計算するために、\par
の1、2、3回微分を計算しないといけない」\par
「たしかに。やや⾯倒ですね」\par
f  ,f\par
c i +1 v  ,f\par
i i\par
f\par
i v\par
i f  =i f(t  ) =i\par
-k(v  ) +i 2 g v\par
i\par
v\par
i+1 v\par
i\par
v  $\simeq$i+1 v  +i\par
=6\par
h(f  + 4f  +f  )i c i +1 (v  のみの式)i\par
i = 0 v\par
$\simeq$1 v\par
0 v\par
0 0\par
v\par
1 i = 1 v  $\simeq$2 v\par
1\par
v\par
1 v\par
2\par
i = 2, 3, 4, ...... v  ,v  ,v  , ......3 4 5\par
v  ,v  ,v  , ......1 2 3\par
f\par
i+1\par
f\par
c f\par
i+1 v\par
i\par
a  ,a  ,a\par
1 2 3 f(t)\par
27\par

% --- PDF page 29 ---
「さらに が変わったら、またその度に微分しないといけない。それはさすがに⾯倒」\par
「ということは、篠澤さんは微分をせずに近似をしたと」\par
「まぁそれに近いイメージかな。具体的には を別の視点から近似して、最終的にシ\par
ンプソン則の式にある を の項まで⼀気に揃える［4］ 」\par
「別の視点？」\par
「さっきのやり⽅は、テイラー展開という式を使って を直接近似した。今から説明\par
するやり⽅は、式というよりかはグラフからそれっぽい値を探すイメージ」\par
「このグラフを⾒てほしい。時刻が のときの速度は にしたよ。このやり⽅では傾きの\par
近さがメイン。 から 、 へ が変わるとき、その傾き（平均変化率）は次のようにな\par
る」\par
「この2つの傾きは、図で⾒るように における の微分係数 に近そう。これを利⽤する\par
んだけど......」\par
続きを⾔いかけたところで、篠澤さんは背伸びして教室の時計を⾒た。数秒悩んだ後、か\par
かとを着け、教壇に持ってきたパソコンをいじりだした。\par
「急にどうしました？」\par
「ちょっとソフトを⽴ち上げてる。今から の近似が、どれくらいの精度かを計算\par
していくんだけど、式変形が⼤変。だから、その式変形を"Maxima"っていうフリーソフ\par
トに全部計算してもらうよ」\par
「そんな便利なソフトがあるんですか」\par
「うん......と、準備ができたね。じゃあ話を戻そう」\par
短くなったチョークをつまみあげ、黒板に体を向ける。\par
f(t)\par
v\par
,v\par
c i +1\par
4f  +c f\par
i+1 h3\par
f\par
,f\par
c i +1\par
t\par
c v\par
c\par
t\par
i t\par
c t\par
i+1 v\par
,\par
2\par
h\par
v\par
-v\par
c i\par
h\par
v\par
-v\par
i+1 i\par
t\par
i v  v˙i\par
f  ,f\par
c i +1\par
28\par

% --- PDF page 30 ---
「まず の近似から。さっきの話から、 の近い値として、 を計算する」\par
「 だから、 は の近い値になるはず。\par
と定めて、 とのズレを で表すことにする」\par
「 をソフトで計算すると......こんな感じ。 は、 の項まで合っていることがわかる\par
ね」\par
「すぐ計算されて便利ですね」\par
「すごい楽になるよ」\par
「次はどうするんですか？」\par
「 が までしか合ってないから、また近い傾きを使って の近似を考えるよ。 は\par
の\par
近い値。⼀⽅で だから、 は時刻 での の微分係数 に近い。\par
が⼗分⼩さければ、 と は似ているはず。だから、 は に近くなるはず」\par
「要は、前に求めた の近似値 を傾きとして使うってこと」\par
「ということは、 で使った を に変える、ということですか」\par
f\par
c v\par
c V\par
1\par
V  =1 v  +i\par
$\cdot$v˙i\par
=2\par
h v  +i f  $\cdot$i\par
2\par
h\par
=v˙ f(t) = -kv +2 g -k(V  ) +1 2 g f\par
c F  =1\par
-k(V  ) +1 2 g f\par
c X\par
1\par
f  =c F +1 X  =1 -k(V  ) +1 2 g +X\par
1\par
X\par
1 F\par
1 h\par
X\par
1 = -  +\par
2\par
h k v kv  -g2 2 i2 ( i2 )\par
6\par
h k v kv  -g 3kv  - 2g3 2 i ( i2 ) ( i2 )\par
-  + ......\{ 48\par
h k kv  -g 15k v  - 15gkv  + 2g4 2 ( i2 ) ( 2 i4 i2 2 ) \}\par
F\par
1 h f\par
c V\par
1 v\par
c\par
=v˙ f(t) = -kv +2 g F\par
1 t\par
c v v ˙c h\par
v˙i\par
v˙c F\par
1\par
v˙i\par
→\par
V  $\simeq$v\par
1 c\par
=f(t  ) = -k(v  ) +gv˙c c c 2\par
$\simeq$  v˙i v˙c\par
$\}$\par
$\}$\par
$\}$\par
$\simeq$v˙i -k(V\par
) +1 2 g =F\par
1\par
f\par
c F\par
1\par
V\par
1\par
v˙i F\par
1\par
29\par

% --- PDF page 31 ---
「そうなるね。 の2つ⽬の近似値 を次のように定めるよ」\par
「同じように 、 を定義しよう」\par
「 を計算すれば......」\par
「 も の項までしか合っていませんね。精度は と同じですか」\par
「でも と を使うと、 を まで合わせられる。まず を計算すると、 の\par
項がなくなる」\par
「 、 だから、これらを代⼊する。 の の係数\par
をまとめて と置いたよ」\par
v\par
c V\par
2\par
V  =2 v  +i F  $\cdot$1\par
2\par
h\par
F\par
2 X\par
2\par
F\par
2\par
f\par
c\par
= -k(V\par
) +g +X\par
2 2 2\par
=F  +X\par
2 2\par
X\par
2\par
X\par
2 =  -\par
2\par
h k v kv  -g2 2 i2 ( i2 )\par
12\par
h k v kv  -g 3kv  - 5g3 2 i ( i2 ) ( i2 )\par
+  + ......48\par
h k kv\par
-g 3k v\par
- 9gkv\par
+ 4g4 2 ( i2 ) ( 2 i4 i2 2 )\par
F\par
2 h F\par
1\par
X\par
1 X\par
2 f\par
c h2 X  +1 X\par
2 h2\par
X  +X\par
1 2 =\par
12\par
h k v kv  -g 3kv  +g3 2 i ( i2 ) ( i2 )\par
-  + ......24\par
h k kv\par
-g 6k v\par
- 3gkv\par
-g4 2 ( i2 ) ( 2 i4 i2 2 )\par
X  =1 f  -c F\par
1 X  =2 f  -c F\par
2 X  +1 X\par
2 h ,h3 4\par
P ,Q\par
→\par
f\par
-F\par
+f\par
-F\par
=Ph +Qh + ......c 1 c 2 3 4\par
f  =  +  h +  h + ......c 2\par
F  +F\par
1 2\par
2\par
P 3\par
2\par
Q 4\par
30\par

% --- PDF page 32 ---
「この式、よく⾒ると という式が、 の\par
の項まで合わせられることを表している」\par
「今までの近似をうまく組み合わせると、 の近似値の精度が上がると」\par
「そうなるね」\par
「じゃあ次は、 を使って を近似してから を計算し、 を組み合わせること\par
で、 の の項まで合わせるということですね」\par
「そうしたいところだけど、ここでプロデューサーに残念なお知らせ。その を計算して\par
も\par
の精度はもう上がらない。このやり⽅だと、 の近似精度は までが限界」\par
「ほう......」\par
「でも、まだ の近似をやり残している。こっちの近似で を定義するよ」\par
「 を計算するよ」\par
「 は まで合いますが、 はまだ残っていますね」\par
「でも、この\par
の項は上⼿に消すことができる。プロデューサー、わかる？」\par
「あんまり思いつかないですね......さっきの流れだと、 と組み合わせるような気が\par
するんですが......」\par
「いいね。その線で合ってる」\par
⾒慣れない数式に⽬が眩む。黒板の下から上へ、右から左へゆっくりと⽬を流すと、ある1\par
つの式に⽬が留まった。\par
「......なるほど。 の の項は、 の の項のちょうど 倍ですね」\par
2\par
F  +F\par
1 2 f\par
c h2\par
f\par
c\par
F\par
2 v\par
c F\par
3 F\par
,F\par
,F\par
1 2 3\par
fc h3\par
F\par
3\par
f\par
c f\par
c h2\par
f\par
i+1 V  ,F  ,X\par
3 3 3\par
V\par
3\par
F\par
3\par
f\par
i+1\par
=v  +F  $\cdot$hi 2\par
= -k(V  ) +g3 2\par
=F  +X\par
3 3\par
X\par
3\par
X\par
3 = -\par
6\par
h k v kv  -g 3kv  +g3 2 i ( i2 ) ( i2 )\par
-  + ......6\par
gh k kv  -g 3kv  +g4 2 ( i2 ) ( i2 )\par
F\par
3 h2 h3\par
h3\par
X  ,X\par
1 2\par
X\par
3 h3 X  +1 X\par
2 h3 -2\par
31\par

% --- PDF page 33 ---
「さすがプロデューサー。それを使うと が消せる。 の の係数を\par
で置き換えて、これが成⽴する」\par
「 の近似式を求めたときと同じようにやってみよう。 、\par
を代⼊する」\par
「この左辺どっかで⾒たことがない？」\par
「篠澤さんが⾔っていた、シンプソン則の計算できない部分ですね」\par
「そう。置き換えてみよう」\par
「この式は、 がシンプソン則の式と の項まで⼀致していることを⽰し\par
ている。前に⾔った通り、シンプソン則は定積分を の精度まで合わせる⽅法。だから、\par
シンプソン則の代わりに を使っても、定積分を の精度まで合わせられ\par
る。これまでの式をまとめて並べるよ」\par
「これも が分かっていれば、 ～ を順番に計算した後に、それらを最後の式に代⼊す\par
ることで、 の近似値が計算できる。だから、 から順に⾬粒の落下速度が計算でき\par
る。このシミュレーションの⽅法を、ルンゲ$\cdot$クッタ法と呼ぶよ」\par
h3 2(X  +1 X  ) +2 X\par
3 h4 R\par
2(X\par
+1 X\par
) +2 X\par
=3 Rh +4 ......\par
f\par
c X  =1 f  -c F\par
1 X  =2 f  -c\par
F\par
,X\par
=2 3 f\par
-i+1 F\par
3\par
→\par
2(f  -F  +f -F  ) +f  -F  = +Rh + ......c 1 c 2 i+1 3\par
4\par
4f  +f  = 2F  + 2F  +F  +Rh + ......c i +1 1 2 3 4\par
6\par
h(f  + 4f  +f  )i c i +1 =  +  + ......6\par
h(f  + 2F  + 2F  +F  )i 1 2 3\par
6\par
Rh5\par
6\par
h(f  +2F  +2F  +F  )i 1 2 3\par
h4\par
h4\par
6\par
h(f  +2F  +2F  +F  )i 1 2 3 h4\par
→\par
→\par
f  = -k(v  ) +gi i 2\par
,\par
$\{$\par
$\{$\par
$\{$V\par
1\par
F\par
1\par
V\par
2\par
=v  +f  $\cdot$\par
i i 2\par
h\par
= -k(V  ) +g1 2\par
=v  +F  $\cdot$\par
i 1 2\par
h $\{$\par
$\{$\par
$\{$F\par
2\par
V3\par
F\par
3\par
= -k(V  ) +g2 2\par
=v  +F  $\cdot$hi 2\par
= -k(V  ) +g3 2\par
v  $\simeq$v  +\par
i+1 i 6\par
h(f  +2F  +2F  +F  )i 1 2 3\par
v\par
i f\par
i F\par
3\par
v\par
i+1 v\par
0\par
32\par

% --- PDF page 34 ---
「この式なら を微分せず、直接 を使うだけで楽ですね。これが篠澤さんの使って\par
た⽅法ですか？」\par
「そう。ただわたしのは の2軸のシミュレーション。求めたのは1軸のシミュレーショ\par
ンだから、もうちょっと複雑になる。でも式の数が増えるだけで、原理も⾒た⽬もだいた\par
い同じ」\par
予測\par
「お疲れさま。これでわたしからの説明はおしまい」\par
第⼀関節くらいまで短くなったチョークが置かれる。彼⼥は⼿をたたき、⽩粉が周りを舞\par
う。\par
「篠澤さんもお疲れさまでした」\par
「バドミントンのシミュレーションの説明が無くてごめんね。ちょっと時間が⾜りなかっ\par
た」\par
「いえ。⾬粒だけでもかなり難しかったので、2軸になったら多分ついていけなくなって\par
いたと思います。でも、篠澤さんのやっていたシミュレーションの概念というか、⼿法の\par
全体像はなんとなく理解しました」\par
「ありがとう。そう⾔ってくれると、説明した甲斐がある」\par
「このルンゲ$\cdot$クッタ法を使うと、⾬粒の落下速度はどうなるんですか」\par
「そこまでプログラムを組むのに時間がかからないから、実際にやってみよっか。直径が\par
mmの⾬粒のとき、 は  /mくらいらしい［5］ 。を  m/s とすると......」\par
そう話しながら、⼩さな⼿でパソコンをカタカタとさせる。しばらくして、画⾯をこちら\par
に向けてきた。\par
「うん。できたよ」\par
「落下し始めてから 秒間の落下速度を計算したよ。最初は急激に速度が上がるけど、 秒\par
もしたら  m/sくらいになって、落下速度は上がらない。落下距離なら、  mもしないう\par
ちに速度の増加は終わることがわかる」\par
f(t) f(t)\par
x,z\par
3 k 0.15 g 9.8 2\par
7 3\par
8 50\par
33\par

% --- PDF page 35 ---
「こんなあっさりと結果が出るんですね」\par
「ルンゲ$\cdot$クッタ法は計算⽅法がシンプルだから、プログラムもそこまで複雑にならな\par
い。だから気楽に結果を出せるよ。ほら50⾏もいってない」\par
篠澤さんがプログラムを⾒せる。\par
「すごい短いですね。これなら⾊々な現象を簡単に計算できそうですね」\par
「数式がわかっていればそうだね。あとはパラメーターの正しい値が必要」\par
「というと？」\par
「今、使っている は、直径  mmの⾬粒専⽤の値。直径が  mmになるとこの値は使えな\par
い。その上、⾬粒が⼤きすぎて歪な形になるから、数式も別になる［6］ 」\par
「ならば、使っているパラメーターの値は、どうやって求めたんですか？」\par
「それは実験から求めている。実験から測定された落下速度によく合う の値が くら\par
いだったということ」\par
「......じゃあ篠澤さんの作ったシャトルコックのシミュレーションは」\par
「もちろん、⾊々なシャトルコックを⾶ばした実験から計算された、それらしい値。今は\par
で設定している［7］ 」\par
「そうなんですね。⾊々な実験で合わせているから、シミュレーションはそれらしい軌道\par
が表⽰されていたと」\par
「そう。でも......」\par
少し⾔葉をためた後、黒板消しを⼿に持ち、上から下へ腕を振り、左から右へ体を動かし\par
⽂字を消していく。続けて彼⼥はこう⾔い続けた。\par
「この軌道は、シャトルコックがラケットから離れた後のやつ。パラメーターとして、打\par
ち出す⾓度と速度を⼿動で⼊れている」\par
「......そのパラメーターはどうやって？」\par
「打ち出す⾓度はラケットの持ち⽅で変わるから、決まった値はない。打ち出す速度は、\par
ラケットの振る速さで変わるから、時と場合による。⼈によっては、あり得ないパラメー\par
ターの組み合わせもあるだろうね」\par
「......そのシャトルコックのシミュレーション、初速度はいくつですか？」\par
「  m/s。時速に直せば、  km」\par
「その軌道、打てるとお思いで？」\par
「まず無理だろうね。腕が砕けちゃうかも」\par
「じゃあ、正しい値はいくつなんですか」\par
k 3 5\par
k 0.15\par
0.26\par
40 144\par
34\par

% --- PDF page 36 ---
「プロデューサー。パラメーターのそれらしい値は、基本的に実験から得られる。でも、\par
わたしはバドミントンの経験があまりない。だから実現できるパラメーターの組み合わせ\par
は、わたしも知らない。だから--」\par
黒板を消し終わり、右から中央へ戻って来る。黒板の前で教壇に両⼿をつき、⽩髪の少⼥\par
は、おそらく今⽇⼀番強い声でこう続けた。\par
「プロデューサー。今から体育館に⾏く」\par
「......はい？」\par
「体育館に⾏く。ラケットを借りて、バドミントンの練習をしながら、パラメーターの組\par
み合わせを⾒つけ出す」\par
「俺が同⾏する必要性は？」\par
「わたしのバドミントンのプレイを動画に撮った後、シミュレーションに合わせて⾓度と\par
初速度を割り出す。動画を取るには2⼈必要。だから、プロデューサーは動画を取る係」\par
「......俺にも予定があるのですが」\par
「1時間くらい前に、午後はとくに⽤事はない、と⾔ったの覚えてる」\par
1時間前の⾃分を反省させたい。\par
「外は⾬が降っています。より良い天気なときの⽅が、気持ちよく練習できると思います\par
よ」\par
ダメ元で反論してみたが、彼⼥は勝利を確信したような笑みを浮かべた。\par
「そうだね。晴れてる⽅が良いかもね。なら......」\par
2⼈で同時に顔を窓へ向ければ、雲の隙間から光が差し込んでいた。\par
「今からがちょうど良さそうだね」\par
どうやら神様は篠澤さんの味⽅をしたらしい。いや、腹をくくれと俺に⾔っているのかも\par
しれない。\par
「ねぇ、プロデューサー」\par
固まっている俺の前で微笑む⼥神は優しく告げる。\par
「コンピューターに計算させた値は、正確な値からズレることがほとんど。パラメーター\par
が多くなると、解く数式が複雑になってズレが⼤きくなることだってある。偉い⼈は、そ\par
のズレをできるだけ⼩さくする⽅法を研究してきた。天気予報はその代表例だと思う」\par
続けて⾔う。\par
「予測精度が \%でなくても、それなりに⾼ければ、その予測を⾒て⼈は⾏動を変える\par
ことができる。今⽇の話も、天気予報を⾒ることで説明する部分としない部分の取捨選択\par
を、最初からでも、途中からでもできる」\par
100\par
35\par

% --- PDF page 37 ---
まさか。\par
「現代の天気予報は1時間ごとの天気まで予測できるみたい。すごいよね。予測精度は\par
\%程度と、割と⾼いんだって」\par
ずっと篠澤さんの⼿のひらの上で俺は踊っていたようだ。\par
70\par
36\par

% --- PDF page 38 ---
おまけ: 落下速度の数式\par
「ところで、⾬粒の落下速度は数式として求まるのですか？」\par
「求まるよ。やってみよっか」\par
＊ ＊ ＊\par
を分数とみなして、次のように式変形する。\par
両辺を不定積分する。左辺は次のようになる。 は積分定数。\par
右辺は （ は積分定数）となるから、次のように、 が求まる。\par
⾃然落下、つまり で とすると、 。よって、 は具体的に次の式になる。\par
dt\par
dv\par
dv =-kv +g2\par
1 dt\par
C\par
L\par
dv$\int$ -kv +g2\par
1 =  dv$\int$\par
(  -  v)(  +  v)g k g k\par
1\par
=  +  dv2 g\par
1 $\int$\par
-  v gk\par
1\par
+  v gk\par
1\par
=  ln    +C\par
2  kg\par
1\par
-\par
v gk\par
+\par
v gk\par
L\par
dt =$\int$ t +C\par
R C\par
R v\par
→\par
→\par
ln    +C  =t +C\par
2  kg\par
1\par
-  v gk\par
+  v gk\par
L R\par
=Ae (A =e )\par
-\par
v gk\par
+\par
v gk 2\par
t kg 2\par
(C\par
-C\par
)kg\par
R L\par
v =\par
k\par
g (Ae + 12  t kg\par
Ae - 12  t kg\par
)\par
t = 0 v = 0 A = 1 v\par
37\par

% --- PDF page 39 ---
＊ ＊ ＊\par
「--って感じ」\par
「うーん。正直なところあまり分かっておりませんが、とりあえず数式はわかるんですね」\par
「そう。シャトルコックに⽐べれば簡単だから」\par
「 の数式が求まるなら、コンピューターで計算するのには意味があるのですか？」\par
「確かに今回のような⾬粒の例だとそこまで困らない」\par
⼿を顎に当て、うーんと、考え始めた。何度⾒ても、この姿は他の誰よりも様になる。\par
「少し状況を変えてみよう。実際の⾬粒の空気抵抗には、落下速度に1乗と2乗の項がある\par
ことが知られている。だから、運動⽅程式に1乗の項、 を加えてみる」\par
「この の式を頑張って解くと、次のようになる」\par
「相当複雑な式になってしまうのですね」\par
「そう。微分⽅程式に を加えただけなんだけどね。もう1つ、極端な例を挙げるね」\par
「初速度が でないとき、この は簡単に表せない。でも、数値計算はできる。 の式が複\par
雑になったり、簡単に表せなかったりする場合に、⼿っ取り早くグラフに表してくれる。\par
これが数値計算のいいところの1つ」\par
v =\par
k\par
g (\par
e + 12  t kg\par
e - 12  t kg\par
)\par
v\par
-mlv\par
m =v˙ - mkv -2 mlv +mg\par
v\par
v =p\par
p\par
ただし、p\par
=\par
+ - (p  -p  e- + k(p  -p  )t- +\par
1 -ek(p  -p  )t- +\par
) ( $\pm$ 2k\par
l $\pm$  l + 4kq2\par
)\par
-mlv\par
m =v˙ - m log(v + 1)\par
0 v v\par
38\par

% --- PDF page 40 ---
参考⽂献\par
［1］ R. Gunn and G. D. Kinzer, "The terminal velocity of fall for water droplets in\par
stagnant air", J. Atmos. Sci. 6 (1949) 243.\par
［2］ tenki.jpチーム. “【⼗種雲形】雲は全部で10種類 ⾒分け⽅を形や⾼さから解説！\par
～下層雲編～”. tenki.jp. 2021-08-01.\par
https://tenki.jp/suppl/tenkijp\_labo/2021/08/01/30532.html, (2025-06-30).\par
［3］ マスオ. “シンプソンの公式の証明と例題”. 学びTimes. 2021-03-07.\par
https://manabitimes.jp/math/766, (2025-06-30).\par
［4］ ⽮部孝, 内海隆⾏, 尾形陽⼀. CIP法ー原⼦から宇宙までを解くマルチスケール解法.\par
森北出版, 2005, 222p.\par
［5］ J. Almedeij, "Drag coefficient of flow around a sphere: Matching\par
asymptotically the wide trend", Powd. Tech. 186 (2008) 218.\par
［6］ B. Bossa and E. Villermaux. "Single-drop fragmentation determines size\par
distribution of raindrops". Nat. Phys. 5 (2009) 697.\par
［7］ 阿部邦昭, 村⽥浩. "シャトルコックの運動". PESJ. 37 (1989) 281.\par
39\par

% --- PDF page 41 ---
26\%の誤算\par
その⽇、広は学園のキッチンを訪れていた。何をしに来たかというと、別に調理をする\par
ためではなく、ただ、友達の花海佑芽と⼀緒に寮に帰りたくて彼⼥を探し、キッチンにい\par
るという話を聞いて、レッスン後の体に鞭打って寄り道しに来たという⼨法だ。広がキッ\par
チンの扉を引いて⼀⾔、\par
「佑芽、⼀緒に」\par
とまで⾔ったその時だった。\par
「うわああああああ！！ もうわかんないよぉおおおお！！！！」\par
貸し切り状態のキッチンに響き渡る⼄⼥の叫びは、放課後の穏やかな静謐を破り、空気と\par
そこらじゅうの⾷器を震わし、そして篠澤広の貧弱の五体へと到達する。刹那、体が宙に\par
浮くような感覚と共に膝から崩れ落ちそうになったが、広はすんでのところでドアの縁を\par
掴み、体を⽀えようと必死に努⼒した。枯れ⽊の枝のような細い腕も、今にもバランスを\par
崩しそうになっている細い⾜も、⼀体そのどこに⼒を⼊れて⽴ち続ければいいのか。⼤丈\par
夫、体が覚えている。そのために私はレッスンを続けてるんだから、ね。と彼⼥が思った\par
かどうかはわからないが、幸いにも今⽇⾏ったのはダンスレッスンではなくボーカルレッ\par
スンだったため、広は⾟うじて体勢を⽴て直し、倒れずに済んだ。やった。やってやっ\par
た。多分彼⼥はそんな気持ちだった。\par
「あ、広ちゃん！ ごめんね、なんかびっくりさせちゃった？」\par
佑芽が駆け寄って来て広に声をかける。\par
「佑芽の声がする。よかった、私の⿎膜、破れてなかった」\par
広は蚊の鳴くような声でそう呟いた。\par
「⼤げさだよ！ ステージの上ではもっと⼤きい⾳するでしょ！？」\par
佑芽にそう⾔われて広は少し⾸を傾げつつも表情を変えず、\par
「⼀理ある」\par
と答えた。そして、\par
「それはそうと、キッチンで何を作ろうとしてたの。もしかして、家庭科の補習？」\par
広がそう尋ねると、佑芽は\par
「違うよ！ これを⾒てもらえばきっとわかると思うんだけど……」\par
と、テーブルの上に載った⼀つのお弁当箱を指さした。それはまさに、おにぎりの形をし\par
たお弁当箱だった。\par
40\par

% --- PDF page 42 ---
「⾒慣れたお弁当箱。けど私はこのお弁当箱、⾯⽩いと思ってる。⾷べ物を⼊れる箱が\par
⾷べ物の形をしてて、マトリョーシカみたい」\par
広は顎に⼿を当ててその「転がり続ける元気の源」を品評し始めた。\par
「おにぎりを⾷べると元気が出る。元気が出れば、きっと佑芽は⾛り続けられる。まる\par
で、ロータリーエンジンみたい。でも佑芽はおにぎりを⾷べてるから、ロータリーエンジ\par
ンを⾷べて⾛ってるってことになる、ね。いや、でもこの場合⾷べてるのはロータリーエ\par
ンジンの中⾝？」\par
そんな意味不明な話をしていると、\par
「広ちゃんなんの話してるの？ よくわかんないけど、あたしね、このお弁当箱になるべ\par
くたくさん⾷べ物を詰めたいんだ。けど、何をどうやって⼊れたらいいか全然わかんなく\par
て」\par
と、佑芽はそう⾔って困った顔をしていた。\par
「なるほど、ね。うーん……もし、本当にたくさん詰め込みたいなら、そもそも容器の\par
形は円柱状の⽅が体積が⼤きくなるんだけど、きっとこのおにぎりの形は佑芽のトレード\par
マークみたいなもののはず、でしょ？」\par
広がそう⾔うと佑芽は\par
「うん！ あたしこのお弁当箱、かわいくて美味しそうで⼤好き！」\par
と答えた。さらに広はもう⼀つ質問した。\par
「佑芽はお姉ちゃんにお弁当を作ってもらえば、ペースト状の⾷べ物でお弁当箱がいっ\par
ぱいに満たされる。それじゃあダメ？」\par
すると佑芽は\par
「最近のお姉ちゃんはね、普通の料理も作るようになったんだ。だから今回はそれはナ\par
シだよ！」\par
と答えた。それを聞いて広は何かを決意したように、\par
「わかった。佑芽は⼤事な友達だから、意志を尊重する。これはいわゆるパッキング問\par
題。さて、と……」\par
と、おもむろにスマホを取り出してアプリでお弁当箱の⼤きさをスキャンし始めた。そし\par
て箱の裏に印字された容量をチェックし、\par
「お弁当箱の⾼さは約7cm、そして容量は900ml……⼤きいね。私なら、300でいい。い\par
や、それだと幼児⽤になっちゃう、ね。ふふ。けどそんなこと⾔ってる間に計算できた」\par
広は以下のような⼿法で、ロータリーエンジンのような形状のお弁当箱の⼀辺の長さを割\par
り出した。\par
41\par

% --- PDF page 43 ---
1.体積の導出\par
この容器は「ルーローの三⾓形断⾯ $\times$ ⾼さ」の柱体とみなせる。柱体の体積は以下の式\par
で表される:\par
ここで、\par
:容積（体積）=900cm$^{3}$\par
:⾼さ=7cm\par
:断⾯積（ルーローの三⾓形の⾯積）\par
よって、⾯積 は以下のように概算できる:\par
2.ルーローの三⾓形の⾯積の導出\par
ルーローの三⾓形の⼀辺の長さを とすると、断⾯積 は以下の式で表される:\par
この定数部分は次のように近似できる:\par
よって断⾯積 の近似式は以下の通りである:\par
V = A $\cdot$ h\par
V\par
h\par
A\par
A\par
A =  =h\par
V\par
$\approx$7\par
900 128.57\par
a A\par
A =  ($\pi$ -2\par
1\par
)a3 2\par
($\pi$ -2\par
1\par
) $\approx$3  (3.1416 -2\par
1 1.732) =  =2\par
1.4096 0.7048\par
A\par
A $\approx$ 0.7048 $\cdot$a2\par
42\par

% --- PDF page 44 ---
3.⼀辺の長さの導出\par
断⾯積の式に1.で求めた を代⼊すると、\par
結論:\par
⼀辺の長さはおよそ13.5cmである。\par
「⼀辺の長さは約13.5cm。ここに、可能な限り多くの⾷材を詰め込めばいい」\par
広はここで⼀つ、驚くべき前提を思いついた。\par
「⾷べ物、全部お団⼦の形にしていい？」\par
すると佑芽は、\par
「えー！？ まあ、広ちゃんがそう⾔うなら……信じるよ！ そうするといっぱい⼊るんだ\par
よね！？」\par
と尋ねてきた。すると広は\par
「⾷べ物にも⾊々な⼤きさや形があるけど、⽴⽅最密充填構造を実現するにあたって、\par
球状であるのが理想的。ミートボールなんかがいいモデルになる、ね」\par
と答えつつ、既にメモ帳とペンを取り出し計算を始めていた。\par
「ミートボール1個の直径が⼤体2.5cm。これを基準にする。つまり、お弁当箱にミート\par
ボールは何個⼊るのか、計算してる……」\par
ここで今、広の頭の中では以下のような考察がなされていた。\par
1.問題設 定のおさらい\par
容器 :断⾯が幅13.5cmのルーローの三⾓形、⾼さ7cmの柱体。\par
球 :直径2.5cm、半径 =1.25cmの完全な球体。\par
容器 に、なるべく球 を詰め込んだ時の詰め込み可能数 を求めたい。\par
a\par
A\par
128.57 $\approx$ 0.7048 $\cdot$a2\par
a $\approx$2 $\approx$0.7048\par
128.57 182.38\par
a $\approx$  $\approx$182.38 13.5\par
A\par
B r\par
A B n\par
43\par

% --- PDF page 45 ---
断⾯となるルーローの三⾓形は、正三⾓形の3頂点を中⼼とする半径13.5cmの円３つの共\par
通部分として表現される。また、球が容器に収まるためには、球の中⼼が容器の境界から\par
半径\par
ぶん内側に存在する必要がある。\par
以下、座標空間上で考える。断⾯は容器の⾼さによって変形しないから、球の中⼼の座標\par
から を⼀つ固定したときのことを考える。\par
正三⾓形の⼀頂点を原点に配置し、もう⼀点を 軸の正の⽅向に、最後の点を 軸の正の\par
⽅向に配置する。また、 軸の正の⽅向に容器が収まるように、座標空間上に容器 を配\par
置する。このとき、正三⾓形の頂点 は以下の座標で表される:\par
この断⾯内に円が収まるには、円の中⼼座標 が以下を満たす必要がある:\par
ここで、  とは、 のうち最⼤の値をあらわすもので、\par
に対し、 をあらわすも\par
の、つまり2点 間の距離をあらわすものとする。 この領域を縮⼩ルーロー三⾓形と呼\par
ぶことにして、計算を簡単にするために、点 たちを以下のように取り直す:\par
ここで、 は縮⼩ルーロー三⾓形の径、つまり であ\par
る。これを代⼊すると、 、 、\par
に位置することが分かる。これは、上で定め\par
た縮⼩ルーロー三⾓形の⼀点が原点に位置するように平⾏移動した図形である。\par
⼀⽅、球の中⼼は ⽅向についても同様の理屈で以下の範囲にある必要がある（点 を取り\par
直したように、 の範囲も取り直している）:\par
r\par
(x, y, z) z\par
x y\par
z A\par
p  , p  , p\par
1 2 3\par
p\par
=1 (0, 0), p\par
=2 (a, 0), p\par
=3\par
,\par
a( 2\par
a\par
2\par
3 )\par
(x, y)\par
max $\parallel$(x,y) -p\par
$\parallel$,  $\parallel$(x,y) -p\par
$\parallel$,  $\parallel$(x,y) -p\par
$\parallel$ $\leq$\{ 1 2 3 \} 12.25\par
max(a, b, c) a, b, c p =\par
(x  , y  ), q =1 1 (x  , y  )2 2 $\parallel$p - q$\parallel$ =  (x  - x  ) + (y - y  )1 2 2 1 2 2p, q\par
p\par
p  =1 (0, 0), p  =2 (a , 0), p  =$\prime$ 3\par
,  a( 2\par
a$\prime$\par
2\par
3 $\prime$)\par
a$\prime$ a =$\prime$ a - r = 13.5 - 1.25 = 12.25\par
p  =1 (0, 0) p  =2 (12.25, 0) p  =3\par
6.125,  12.25 $\times$  $\approx$( 2\par
3 ) (6.125,  10.608)\par
z p\par
z\par
44\par

% --- PDF page 46 ---
⽬的: 上記の条件に従い、容器 に球 を詰め込む際の球の中⼼の座標を求める。詰め込\par
み⽅としては、結晶構造を参考にして、単純⽴⽅(sc)格⼦、六⽅最密充填(hcp)構造、⾯⼼\par
⽴⽅(fcc)格⼦の配置になるように球を詰め込み、詰め込める最⼤数をシミュレーションす\par
る。\par
2.1 SC（単純⽴⽅格⼦）\par
SCは最も原始的な格⼦で、隣り合う球の中⼼が⽴⽅体になるように球を敷き詰める。 球\par
の直径が2.5cmであるから、球の中⼼の座標を\par
から各⽅向に2.5ずつ\par
ずらして容器内に配置できるか確認する。\par
軸⽅向には同じ断⾯をしているので、⾼さと断⾯を分けて考える。 ⾼さは7cm、球の直\par
径が2.5cmであるから、2層積むことが出来る。 球の中⼼を通る断⾯を考え、その断⾯に\par
円がどれだけ敷き詰められるかを考える。\par
SCは平⾯上での円の中⼼が正⽅形になるように敷き詰める構造であることを踏まえると、\par
直径が2.5cmであるから、\par
から 軸⽅向に ずつ、 軸⽅向に ずつずらしながら円を\par
配置し、その円の中⼼が縮⼩ルーロー三⾓形の内部に位置していればよい。判定には、円\par
の中⼼と縮⼩ルーロー三⾓形を構成する3点それぞれとの距離がいずれも径より⼩さいか\par
どうかを⽤いる。以下のPythonスクリプトを⽤いてシミュレーションすると、15個配置\par
されることが分かる\par
。\par
import math\par
\# 3点の座標を与える\par
k = 12.25 \# 縮小ルーロー三角形の径\par
d = 2.5   \# 球の直径\par
r3 = math.sqrt(3)\par
p\_arr = [[0,0], [0,k], [k/2, r3]]\par
\# 縮小ルーロー三角形の内部の点を数える\par
count = 0\par
\# 単純立方格子は立方体の頂点に球の中心が来る構造\par
\# 球の中心をずらしていく\par
for j in range(7): \# y軸に沿って動かす。余裕をもって0から6までシミュレーション\par
する。\par
for i in range(7): \# x軸に沿って動かす。余裕をもって0から6までシミュレーシ\par
ョンする。\par
\# 球の中心が縮小ルーロー三角形の内部に位置するか判定\par
0 $\leq$z $\leq$ 7 - 2r = 7 - 2.5 = 4.5\par
A B\par
(x, y, z) = (0, 0, 0)\par
z\par
p\par
1 x 2r y 2r\par
1\par
45\par

% --- PDF page 47 ---
is\_contained = True\par
\# 各頂点との距離が径よりも小さければ内部に存在する\par
for p in p\_arr:\par
if math.sqrt( (p[0] - i*d)**2 + (p[1] - j*d)**2 ) > k:\par
is\_contained = False\par
if is\_contained:\par
count += 1\par
\# 総数を出力\par
print(count)\par
2層詰められるので、合計で30個となる。\par
3.2 HCP（六⽅最密充填）\par
HCP は “最密充填” のひとつである。HCPの場合、層によって配置がズレるのでそれぞれ\par
の層での球の中⼼の座標を考える必要がある。⼀番下の層では中⼼どうしが正三⾓形を構\par
成するように、真ん中の層では下の層の球の隙間を埋めるようにして、次の層ではまた最\par
下層と同じように詰め込んでいく配置⽅法である。各層の詰め込み⽅を考察すると、正三\par
⾓形の性質から、隣り合う球の中⼼の座標は\par
軸⽅向には直径分ズレるが、 軸⽅向には\par
直径の 倍ズレる。 また、 軸⽅向を考えると正四⾯体ができるので、2層⽬球は1層⽬\par
と⽐べて 軸⽅向に直径の 倍、 軸⽅向に 倍、 軸⽅向には直径の 倍ズレる。これ\par
らを踏まえてSCと同様に以下のPythonスクリプトを⽤いてシミュレーションを⾏うと、\par
61個詰められることがわかる\par
。\par
import math\par
\# 3点の座標を与える\par
a = 12.25 \# 縮小ルーロー三角形の径\par
d = 2.5   \# 球の直径\par
r6 = math.sqrt(6)\par
r3 = math.sqrt(3)\par
p\_arr = [[0,0], [0,a], [a/2, r3]]\par
\# 縮小ルーロー三角形の内部の点を数える\par
count = 0\par
\# 六方最密充填(HCP)はハニカム構造の隙間を埋め尽くして積み上げるような構造\par
\# 球の中心をずらしていく\par
for k in range(7): \# z軸に沿って動かす。余裕をもって0から6までシミュレーション\par
する。\par
\# 高さがはみ出したら中断\par
x y\par
2\par
3 z\par
x\par
2\par
1 y\par
6\par
3 z\par
3\par
6\par
1\par
46\par

% --- PDF page 48 ---
if k*r6/3*d > 7-2.5:\par
break\par
for j in range(7): \# y軸に沿って動かす。余裕をもって0から6までシミュレーシ\par
ョンする。\par
for i in range(7): \# x軸に沿って動かす。余裕をもって0から6までシミュレー\par
ションする。\par
\# 球の中心が縮小ルーロー三角形の内部に位置するか判定\par
is\_contained = True\par
\# 各頂点との距離が径よりも小さければ内部に存在する\par
\# 構造上、z軸偶数と奇数の層、y軸偶数と奇数の層で座標がすれるので補完する\par
for p in p\_arr:\par
if math.sqrt( (p[0] - (i*d + (d/2)*(k\%2) - (d/2)*(j\%2) ) ) **2 +\par
(p[1] - ( (r3/2)*j*d + (r3/6)*(k\%2) ) )**2 ) > a:\par
is\_contained = False\par
if is\_contained:\par
count += 1\par
\# 総数を出力\par
print(count)\par
3.3 FCC（⾯⼼⽴⽅格⼦）\par
FCCの場合、HCPと考え⽅は同じだが、3層⽬の積み⽅が異なる。HCPの際は1層⽬と同\par
じ位置に詰め込んだが、1層⽬とも2層⽬とも異なる位置であって、最密に詰め込める隙間\par
が存在する。この時その隙間は、球の中⼼の位置を上から⾒た時、1層⽬の正三⾓形の部\par
分に着⽬すると、2層⽬はその正三⾓形の重⼼に、3層⽬は重⼼と対称な点に位置するよう\par
な形となる。つまり3層⽬は1層⽬と⽐べて\par
軸⽅向にズレず、 軸⽅向に 倍ズレた並び\par
⽅をする。こちらもSC、HCPと同様にPythonスクリプトを⽤いて検証を⾏う。\par
import math\par
\# 3点の座標を与える\par
a = 12.25 \# 縮小ルーロー三角形の径\par
d = 2.5   \# 球の直径\par
r6 = math.sqrt(6)\par
r3 = math.sqrt(3)\par
p\_arr = [[0,0], [0,a], [a/2, r3]]\par
\# 縮小ルーロー三角形の内部の点を数える\par
count = 0\par
\# 面心立方格子(FCC)はハニカム構造の隙間を埋め尽くして積み上げるような構造\par
x y\par
3\par
3\par
47\par

% --- PDF page 49 ---
\# 球の中心をずらしていく\par
for k in range(7): \# z軸に沿って動かす。余裕をもって0から6までシミュレーション\par
する。\par
\# 高さがはみ出したら中断\par
if k*r6/3*d > 7-2.5:\par
break\par
for j in range(7): \# y軸に沿って動かす。余裕をもって0から6までシミュレーシ\par
ョンする。\par
for i in range(7): \# x軸に沿って動かす。余裕をもって0から6までシミュレー\par
ションする。\par
\# 球の中心が縮小ルーロー三角形の内部に位置するか判定\par
is\_contained = True\par
\# 各頂点との距離が径よりも小さければ内部に存在する\par
\# 構造上、z軸の層、y軸偶数と奇数の層で座標がすれるので補完する\par
for p in p\_arr:\par
\#層の位置によってずれを調整する 2層目まではHCPを流用\par
if k\%3 != 2:\par
if math.sqrt( (p[0] - (i*d + (d/2)*(k\%2) - (d/2)*(j\%2) ) ) **2\par
+ (p[1] - ( (r3/2)*j*d - (r3/6)*(k\%2) ) )**2 ) > a:\par
is\_contained = False\par
else:\par
if math.sqrt( (p[0] - (i*d + (d/2)*(j\%2) ) ) **2 + (p[1] - (\par
(r3/2)*j*d - (r3/3) ) )**2 ) > a:\par
is\_contained = False\par
if is\_contained:\par
count += 1\par
\# 総数を出力\par
print(count)\par
結果、57個詰められることが分かる 。\par
4.結果まとめ\par
配置⽅式 球数 層数 1層あたり 密度の印象\par
SC 30 2 15 低い\par
HCP 61 3 約20.3 ⾼い\par
FCC 57 3 19 ⾼い\par
1\par
48\par

% --- PDF page 50 ---
5.考察と結論\par
理論密度はFCCとHCPが同じ だが、断⾯形状へのフィット性によって実際の\par
球数に差が出る。\par
今回の断⾯（ルーローの三⾓形）では、HCPとFCCの切り出し形がいずれもよく適合し、\par
HCPは最⼤の61個が配置可能だった。\par
FCCは理論上は同密度だが、今回のシミュレーション条件（ を起点に敷き\par
詰めた場合）ではわずかにHCPに劣る形となった。\par
SCはもっとも単純だが隙間が多く、30個が限界である。\par
結論: 本ケースでは概ねHCPが最適配置。\par
「できた。61個」\par
広はそう結論を出した。佑芽は\par
「わぁーっ！！ すごいね広ちゃん！」\par
と⽬を輝かせた。広もこれにはさすがに得意げな顔をしている。\par
「広ちゃんの⾔ってることは全然わかんなかったけど、とにかく広ちゃんの⾔うとおり\par
に詰めれば61個⼊るんだよね！ それにしても、あぁー、お弁当箱に61個もミートボールが\par
詰まってるところを想像したら、お腹減ってきたなぁー。試してみるのが今から楽しみだ\par
よ～」\par
と⼤喜び。その様⼦を⾒て広もまた、穏やかな笑みを浮かべていた。⼤好きな友達の、役\par
に⽴つことができた。そのことが何よりも嬉しかったからだ。しかし、そこに嵐が訪れ\par
る。\par
「佑芽！ こんなところにいたのね。廊下まで声が漏れてるわよ」\par
そこに現れたのは花海咲季。⼾の前で両の⼿を腰に当て仁王⽴ちしているその姿は、ステ\par
ージの幕開けを思わせる堂々たる佇まいであった。\par
「あ、お姉ちゃん！！ 今ね、広ちゃんがお弁当箱にどうやって⾷べ物を詰めたらたくさ\par
ん⼊るか相談してたんだ。お団⼦の形にすれば61個も⼊るんだって。すごいよね！」\par
佑芽はまるで⾃分で発⾒したことかのように誇らしげに、元気いっぱいにその事実を姉に\par
伝えた。広もまた、佑芽のその姿を⾒て⾃分のことのように嬉しい気持ちになっていた。\par
だが⼀⾔、咲季の⾔葉が状況を⼀変させる。\par
「球状？ お弁当箱にたくさん⼊れたいなら、前みたいに全部ペースト状にしたら隙間な\par
くびっちり⼊るんじゃないの？」\par
姉の⾔葉に対してすかさず佑芽が⾔い返す。\par
「そうだけど！！ お姉ちゃん最近はそういうのじゃない料理も作ってるよね。あたしも\par
0.74048 …\par
(x, y) = (0, 0)\par
49\par

% --- PDF page 51 ---
噛んで⾷べるタイプのお弁当が作りたいから！」\par
すると咲季は少し⾸を傾げ、\par
「何よ、それならそのお弁当箱⼀杯にお⽶を詰め込んだらいいじゃない。まあ、栄養バ\par
ランス的には問題がありそうだけど……佑芽はおにぎり⼤好きなんだから、お弁当丸ごと\par
おにぎりにしちゃえばいいと思うわ」\par
と⾔い放った。それを聞いて広は突如⼼が折れ膝から崩れ落ちた。\par
「広ちゃん！？？！？！？」\par
佑芽はとっさに倒れ込んできた広を⽀え、そっと膝に乗せるようにその精巧な西洋⼈形の\par
ような華奢な体を抱き寄せた。\par
「ふふ……どうして私、そんな簡単なことに気づかなかったんだろう。やっぱり、佑芽\par
のお姉ちゃんは凄腕のシェフ、だね」\par
広は中空に⼿を伸ばしながら、誰に⾔うでもなくそう呟いた。琥珀⾊の瞳に揺らめく輝き\par
が雫となり、⽬尻を伝い静かに流れ落ちていく。そして広はそっと⽬を閉じた。\par
「⼤好きな友達の役に⽴てたと思ってた。それがたまらなく、嬉しかったはず、なの\par
に。ああ、本当に、ままならない、ね……」\par
広はそう⾔うと、⼒なく意識を失った。\par
「広ちゃーーーーーーん！！！！」\par
佑芽の慟哭は、開いたままのキッチンのドアを突き抜け、放課後の⽇差しを受け暖かな光\par
をたたえる学園の廊下をぶち抜き、どこまでも響き渡っていった。なお、広はその後いつ\par
も通り保健室に運ばれ、10分くらい寝た後、熱⼼に付き添ってくれた花海姉妹と⼀緒に普\par
通に寮に帰って⾏った。果たして明⽇の佑芽のお弁当箱には64個のお団⼦が詰まっている\par
か、それともお⽶がぎっしり詰まっているのか。それは未だ、誰も知らない。\par
脚注\par
1. から敷き詰めた時の結果。起点をずらすことで増減する可能性がある(x, y) = (0, 0)\par
50\par

% --- PDF page 52 ---
すべてがエコーすること、もしくはモンテカ\par
ルロ法と焼きなまし法について\par
イカ爆弾∕⼭船\par
買い被られるひととは天才である。 「天才」を前にしたとき、ひとは黄⾊を⽩に誤解し\par
て、紺を黒だと⾔い間違えてしまう。あんまり明るかったり暗かったりして、その微細な\par
違いには興味が無い事が多いからだ。無論、普通に暮らしている分にはそれでも⼤した問\par
題は起こらない。黒と⾔って取り出した服が⽩かったら問題だろうが、紺⾊ならほぼ同⾊\par
だ。ほぼ同⾊だと了解されるなら、それはほぼ同義なのだ。\par
コンピュータもまた天才のひとつに数えられる。その計算速度は普通の⼈間からすれば\par
常軌を逸している。どんなにめんどくさい計算を放り込んだとて⼀瞬で回答が表⽰される\par
わけだ、やはり普通に計算している分にはそんな⾯倒なことは起こらない。\par
「……そんな認識だったの？」\par
「⾯⽬ない……と⾔うのも適正ではないように思いますがね」\par
「わたしにできないことがすぐそこにあるのと同じくらい、パソコンも計算能⼒は⾼が知\par
れてる、よ」\par
「と⾔われましても……」\par
まだ午後になった実感もないような⾼い陽の下、条件付き天才少⼥とそのプロデューサー\par
は部屋で雑談に興じていた。いた、というのはここで過去形である。紆余曲折は省略する\par
が、このプロデューサーという男がうっかりコンピュータの計算量の話に⾜を突っ込み、\par
それを篠澤が咎めたというだけの話だ。うだうだここを広げたって読者のオタクのみなさ\par
んも困るだろうからさっさと本題に⼊ろう。\par
「……たとえば、密度汎関数法による第⼀原理計算のことを考えたい。この詳細は知らな\par
くてもいいけど、要点は、原⼦配置によってポテンシャルが決定され、そのポテンシャル\par
を踏まえて2変数を決定しなければならないということ。\par
Schrödinger⽅程式におけるエ\par
し ゅ れ ー で ぃ ん が ー\par
ネルギー固有値と波動関数にそれぞれ対応する。密度汎関数法では、これをイテレーショ\par
ンにより解く。つまり、初期値を与えて計算して、出⼒を再度⼊⼒に与えて計算を進め、\par
期待される誤差が⼀定以下になったら打ち⽌めとする、というようなやり⽅。実際にパッ\par
ケージを使ってみればわかるけど、ごく簡単な例でもそこそこ計算に時間がかかる。たと\par
えば、\par
⾯⼼⽴⽅格⼦の純⾦について求めてみると、わたしが前に使ってたパソコンでは5\par
Face-Centered Cubic, FCC\par
分程度だった。\par
プロデューサー、これは実験で予め明らかになっている固定的な配置をもとに原⼦位置\par
を変えずに計算できたからこの程度の時間で済んだ、ってことはわかる、よね。なら、も\par
し計算したい系の構造が実験で明らかになっていなかったり、あるいは実験的最適配置か\par
51\par

% --- PDF page 53 ---
ら⼤きく変形させた配置で原⼦配置の最適化を⾏いたいとき、まともに使える計算量はど\par
のくらいか、わかる？」\par
「……Landauの記号を⽤いれば良いんですよね？ そうなると……ちょっと待って下さ\par
い、そもそも  とするときの  って今回の場合どれにあたるんですか？」\par
「プロデューサー、もしかしてSchrödinger⽅程式は勉強したことない？」\par
「理系ではないんですよ、俺は。プロデューサー科の⼤学⽣ですったら」\par
「⾼校までの課程に出てこないの？ 化学科のカリキュラムでは1年からその応⽤を扱\par
う、よ。だから、みんな知ってると思う。……簡単に説明すると、Schrödinger⽅程式は\par
電⼦と原⼦核によっていろいろ決定するもの。特に⽔素原⼦のものが2体問題として解け\par
るから⼊⾨として有名だと思う。とにかく、今回使ってみるのは……さっき例に出したか\par
ら、Auにしてみよう。こういう重い原⼦なら、電⼦数が⽀配的になる。とはいえ内殻電\par
⼦までいちいち全部使ってると計算量がかさむばっかりでうれしくないから、⼀部電⼦は\par
原⼦核にまとめる擬ポテンシャル法を採⽤する。あとは平⾯波基底とかいろいろ絡んでく\par
るけど、要するに、単位格⼦あたり、すなわち4原⼦あたりだいたい100電⼦ほどがあるこ\par
とがわかれば⼗分。\par
あとは密度汎関数法の基礎としてのKohn-Sham⽅程式の性質から、単電⼦近似によっ\par
こ ー ん $\cdot$ し ゃ む\par
て1電⼦ごとに2体問題が解かれることになる。つまり、ポテンシャルを得る段階でざっと\par
を得る。実際にはここにk点を与えて逆空間の分割を考えたりとかカットオフエネ\par
ルギーを与えたりしてもっと計算量がかさむけど、⾯倒だし本筋ではないからやっぱり省\par
略。\par
……むしろ、量⼦計算を考えるよりも古典的ポテンシャルを考えたほうが良いかも。ポ\par
テンシャル計算を2体間のそれの総和として考えるなら、\par
回の計算が求められる。そ\par
のポテンシャルの中で、電⼦を1つ動かしてみることと1つ追加してみることには⼤した差\par
が無いから、後者を考える。ここから最適配置のためのアルゴリズムの計算量は、悪くて\par
も\par
くらいが好ましい、ね。積として計算量が  程度になる。これは\par
程度ならその4乗で  回程度、つまり0.1秒程度で済む計算量だけど、例えば単位格\par
⼦を2x2x2にした拡⼤格⼦のことを考えてみれば、だいたい10倍の電⼦数になるわけだか\par
す ー ぱ ー せ る\par
ら、これは  回程度の計算回数になって⼀気に所要時間は1000秒ほどになっちゃう。\par
1000秒なら研究向けの計算量としては耐えられる。でも、もっと⼤きくすることは考えら\par
れない。⼩さな系ならかろうじて有効な最適化アルゴリズムの計算量で\par
程度、っ\par
てこと。\par
プロデューサー、正 準 集 団のことは知ってる、よね。普通、正準集団として取り扱わ\par
かのにかる$\cdot$あんさんぶる\par
れる粒⼦の配置を考えるなら、それは正 準分布に従わなくちゃいけない。だから、 」\par
かのにかる\par
「……すみません、前に聞いたとは思うのですが……何でしたっけ、それ？」\par
篠澤の⽬はステップの草原みたいになった。⼀呼吸だけ。それから顎に⼿が付けられた。\par
O(n) n\par
O(n )2\par
C\par
n 2\par
O(n )2 O(n )4 n =\par
100 108\par
1012\par
O(n )2\par
52\par

% --- PDF page 54 ---
「もしかして、プロデューサーが⾼校を卒業してから課程が変わった？」\par
「そもそも⾼校では扱いませんよ、その……それは」\par
「そんなはずはない。授業中に暇だったから教科書は全部読んだ。ちゃんと理想気体の状\par
態⽅程式の話とその導出も乗ってた。だから、扱うはず。この分で⾏けば、3年のはじめ\par
のあたりに配当されると思う。理想気体の状態⽅程式を導出するなら⼩ 正 準 集 団の話\par
みくろかのにかる$\cdot$あんさんぶる\par
になると思うけど、本質的に⼤差は無いから問題は無い、はず。 」\par
「索引は⾒ました？ 明⽰されていないものは扱われませんよ。発展的内容に触れる教員\par
ならあるいは別かもしれませんが」\par
「……『発展的』っていうくくりだと、多分時間が⾜りなくなる。まあいいか。本題に戻\par
ると、割と⾃然な粒⼦集団の構成がいくつかあって、そのうちのひとつが正準集団にな\par
る、ってことだけ覚えてほしい。逆に、粒⼦がこの集団の振る舞いに制約される、ってこ\par
と。\par
話は逸れたけど、まず、経緯としてモンテカルロ法の話をする。これは粒⼦の統計的振\par
る舞いについてシミュレーションしたいというときに、各粒⼦の状態推移を確率過程で表\par
せれば満⾜するんじゃないか、っていうもの。もし各粒⼦がどのような確率で特定の状態\par
に推移するかが予めわかっているなら、あとはそれを計算機を使って追うだけですべての\par
サンプリングされた粒⼦の状態が追える。そうしたら、あとは総和を取ったり平均を取っ\par
たりすれば物理量が得られる。定義上、ここからはほとんど任意の物理量が得られること\par
になる。例えば、エントロピー、圧⼒、\par
Gibbsエネルギーとか。もちろん、これらがわか\par
ぎ ぶ ず\par
るなら、これらから組み⽴てられるすべての物理量も計算できる。\par
ところが、粒⼦の確率過程なんかはわからない。⽔素原⼦のSchrödinger⽅程式なら解\par
析的に解けるし、それだったらたしかに確率過程も出せるけど、集団を取り扱いたいなら\par
絶対的に3体以上の問題になるし、粒⼦数\par
に対して  次元になるSchrödinger⽅程式\par
は数値的に追うだけでも難しくなっちゃう。さっきの例なら、ざっと300次元は必要にな\par
る。第⼀原理計算でもそんなことはやらない。\par
じゃあどうすればいいかっていうと、逆に考える。粒⼦集団の振る舞いの制約が満たさ\par
れれば満⾜ということなら、粒⼦をいくつか、もしくは1つ選んで、その移動先について\par
制約を満たすかどうか調べて移動するかそのままか、そういうことを判定する。具体的に\par
は、移動後のポテンシャルエネルギーで考えるのが簡便。ポテンシャルは原理的には\par
Schrödinger⽅程式で、実際的には簡略化の結果のKohn-Sham⽅程式によって解かれる\par
けど、どちらにせよ、移動後の位置によって再度計算される必要がある。移動に関わず特\par
性の静的なポテンシャルを採⽤できるわけじゃなくて計算し直さなきゃいけない、という\par
ことが重要。\par
それで、やっと本題。ここで移動を許可するかどうかを考えるとき、Metropolis法に準\par
め と ろ ぽ り す\par
拠することが多い。Metropolis法っていうのは、粒⼦の遷移について……通常は空間内の\par
移動について、\par
n 3n\par
53\par

% --- PDF page 55 ---
1. 初期状態を与え、系の総エネルギー  とポテンシャルエネルギー  を計算し、\par
とする。\par
2. 粒⼦をひとつ選んで、移動する。\par
3. 移動後のポテンシャルエネルギーを計算し、\par
1. 移動前とのポテンシャルエネルギー差  が負ならば（つまり移動後のほうが低\par
いエネルギーになるならば）そのまま移動する。但し、  を新たな\par
とする（つまり差分を代⼊する） 。\par
2. 移動前とのポテンシャルエネルギー差  が正ならば  のときのみ移\par
動し、同様に  を処理する。判定に漏れた場合は何も移動しないままそのステ\par
ップを完了する。\par
4. 頭に戻る。\par
というような推移をする。ここでの『特定の確率』にうまい値をあてれば、その推移を任\par
意の集団のものにすることができる。たとえば、正準集団にしたいなら、その確率  は、\par
Boltzmann定数  および温度  によって、\par
ぼ る つ ま ん $\cdot$ こ ん す た ん と $\cdot$ $\cdot$\par
となる。操作を⾒ればわかると思うけど、この  は正準集団に従わせたいときのトリッ\par
ク。\par
プロデューサー、この操作からわかることはある？」\par
「低い⽅へは確実に移動できて、⾼い⽅へは留保が付く、ということですね。……最⼩値\par
を求めるアルゴリズムなら、始めからより⾼い⽅へ移動しなくてもよいのでは？」\par
「そういうアルゴリズムもある、よ。⼭登り法とか。でも、考えてみて。極⼩値が複数あ\par
ったら、どうする？ 原⼦配列の表⾯のような周期的ポテンシャルが期待できるところで\par
そういう⾵な低い⽅へだけ向かうアルゴリズムがあると、困ることがありそう。だから、\par
その壁も超えられるといい、ね」\par
「⽳を開けるか、抜け道を探すか……ってトンネル効果の話だったんですか？」\par
「違うと思う。とにかく、さっきの例では正準集団に従わせたいときのためのトリックが\par
あって、⼀般の最⼩値探索に⽤いるには煩雑。遷移に成功する確率\par
について、\par
E U\par
K  :=D E - U\par
$\Delta$U\par
K  -D $\Delta$U K\par
D\par
$\Delta$U $\Delta$U < K\par
D\par
K\par
D\par
p\par
k\par
B T\par
p(K  ) $\propto$D exp -  ( k  TB\par
K\par
D )\par
K\par
D\par
p\par
p = exp -  ( T\par
$\Delta$U )\par
54\par

% --- PDF page 56 ---
としてみても⽬的は失われない。実際、これはさっきの式について分母に定数を⾜して\par
としたものに相当するから、  を適当に操作することで同じ結果が得られる。よ\par
かったね」\par
篠澤は諳んじるように式を紙に書いていく。書かれた式にプロデューサーの視線は刺さり\par
にいった。的に刺さったダーツのようだ。そのダーツが実はストローであって、じわじわ\par
と理解と直観とをプロデューサーの頭に輸送するのだが、それはあまりに遅いので飽き飽\par
きしてしまう。春夏秋、また。\par
「なんとなくは……わかったことにして良いと思います。具体例などがあるともっと良い\par
のですが」\par
「そう⾔うと思って、論⽂は⽤意してある。むしろ最初からこれで説明すれば良かったか\par
も。さっきの説明では\par
は所与であるか何らかの⽅法により計算したか、ということ\par
を仮定した。けど、実際には何らかの計算⽅法を与える必要がある。せっかくだから第⼀\par
原理計算でエネルギーを与える例を使うと、Nudged Elastic Band, NEB法がある。\par
な っ じ ど $\cdot$ え ら す て ぃ っ く $\cdot$ ば ん ど\par
単に最⼩値を探索するだけでなくて、最⼩値へ遷移する経路まで知りたいことがある。\par
拡散のすべてがそう。拡散は材料表⾯の吸着や固溶などでも⾮常に重要だから、それを適\par
切に計算できるようにしたい。ところが、式を⾒てもらえれば分かると思うけど、焼きな\par
まし法の単純な適⽤ではエネルギー\par
の経路積分値が最⼩となる経路が⾒つけられると\par
は限らない、どころかより⼀般にはほとんど確実に無駄のある経路が通られることにな\par
る。これをこーなーかってぃんぐって⾔う。\par
特に周期的なポテンシャルを考えると、エネルギー最⼩値を持つ点は無限個周期的に並\par
んでいることがわかる。そのそれぞれを中⼼にすり鉢状のエネルギー曲⾯を想定できるか\par
らそうすると、最⼩値同⼠を結ぶ直線の中点に鞍点ができるような曲⾯になるようになめ\par
らかにできる。極⼩値を1点求めたいなら、そこから式の温度をゼロにしたもので遷移し\par
続ければいい。問題は経路の⽅。\par
経路上を何点かに分割すると、その点の間の⼒……まあ⼒はポテンシャルエネルギー差\par
と本質的に同じだからエネルギーの⽅で考えておくと、それが等しくなるようにすること\par
ができる。経路上のサンプリングされた点を動かしながら、すべての点間エネルギー差が\par
等しくなるように収束させれば、ひとまず経路積分値の最⼩化と同じことができる。これ\par
がえらすてぃっく$\cdot$ばんど法。点間のバネにかかる⼒を等しくするイメージ。ちょっと想\par
像してみればわかるけど、これには極⼩値が極めて多く存在する。斜⾯にボールを転がし\par
ていくと同じ始点と終点を通る無数の経路があるのと同じ。\par
だからといって直接ここに焼きなまし法を適⽤するのはまだ筋が悪い。遷移⽅向に対し\par
て垂直な⽅向のエネルギー差も考慮して、合⼒を最⼩化する、としても本質を失わない。\par
だからそうする。そうすると、遷移の経路が発⾒できる。経路は遷移の峠越えを含むか\par
ら、物質の活性化エネルギーが発⾒できることになる、よ」\par
「活性化エネルギーが……わかると、何が嬉しいんでしたっけ」\par
k  =B 1 T\par
$\Delta$U\par
U\par
55\par

% --- PDF page 57 ---
篠澤はかぶりを振って肩をすくめた。なお、その仕⽅は本場アメリカの仰々しさである。\par
「いくらでもある。反応が化学のすべてだから、ね」\par
参考⽂献\par
1. 神⼭新⼀, 佐藤明: 分⼦シミュレーション講座 モンテカルロ$\cdot$シミュレーション（新装\par
版）, 朝倉書店, (2020)\par
2. Maragakis, P., et al: Adaptive nudged elastic band approach for transition state\par
calculation, The Journal of chemical physics 117.10, (2002), pp. 4651-4658\par
56\par

% --- PDF page 58 ---
\noindent\textit{[第 58 页没有可抽取文字层。]}\par

% --- PDF page 59 ---
\noindent\textit{[第 59 页没有可抽取文字层。]}\par

% --- PDF page 60 ---
\noindent\textit{[第 60 页没有可抽取文字层。]}\par

% --- PDF page 61 ---
\noindent\textit{[第 61 页没有可抽取文字层。]}\par

% --- PDF page 62 ---
寄稿者陣あとがき\par
篠澤さん……汎関数微分について教えて下さい……\par
さておき、材料はいいぞ。異種原子の接点はすごいぞ。Pauliですら「表面は悪魔が\par
作った」って言ってるけど現代計算機の力で叩けばわりとどうとでもなる！\par
すみません嘘です、どうとでもはなりません　篠澤さん……すべての計算をO(nlogn)\par
ぐらいにできるウルトラつよつよアルゴリズムを教えて下さい……\par
イカ爆弾 / 山船(@Esqrt\textbf{[缺字]})\par
できるだけ簡易な説明でポアソン分布の話をしたかったのに、気が付いたら篠澤さん\par
がアクセルベタ踏みしてて終わった\par
tokio(@dtokyo\textbf{[缺字]}\textbf{[缺字]}\textbf{[缺字]}\textbf{[缺字]}\textbf{[缺字]}\textbf{[缺字]})\par
ルンゲ$\cdot$クッタ法って、有名なくせに証明はすっ飛ばされがちな、可愛そうな子なん\par
ですよね。Webのどこを探しても、ゴリゴリの計算による証明。エレガントな証明\par
どっかにあるでしょって、思ってた時期があったんですよ。やってやらと意気込んで\par
いたら、ゴリゴリの数式をさっさと捌く篠澤さんと、それに必死に食らいつくPの妄\par
想ができてしまいました。拙い文章ですが、楽しんでいただければ幸いです。\par
うにかじき(@unikajiki\textbf{[缺字]}\textbf{[缺字]})\par
数式を読むのって楽しい……と思える人は少数派だと思いますが、数式によって世の\par
中の悩み事が解決されるのは面白くはないでしょうか。様々な篠澤さんとのやり取り\par
から奥深さの片鱗を味わっていただければと思います。\par
betty(@betty\textbf{[缺字]}\textbf{[缺字]}\textbf{[缺字]})\par
転がり続ける元気の源、と言われたら転がり続けてるんだからロータリーやろなぁ、\par
と思いますし、そうするとルーローの三角形になると思うのですが、実際の原作ゲー\par
ム内のアイテムは普通に三角形だったと思います。てへ。あとオチは「ペースト状な\par
らいっぱい入る」って話を想定していたのですが、書いてる最中に原作ゲーム内で咲\par
季ちゃんが普通にお料理上手になってしまったので内容を変更してお届けする運びと\par
なりました。ありがとうございました。\par
Najiko(＠najiko\textbf{[缺字]}\textbf{[缺字]})\par
61\par

% --- PDF page 63 ---
このたびは本合同誌を手に取っていただき、ありがとうございます。\par
「篠澤広というキャラクターの新たな解釈を模索する」という目標を掲げた本誌『完\par
全理解！篠澤広の機械学習$\cdot$数理最適化ゼミ合同～1限目～』は、寄稿者5名と主宰の\par
計6名による合同誌となりました。\par
本誌は昨年末、冬コミの話題で盛り上がっていたときにふと「わたしも即売会に行っ\par
てみたい」とつぶやいたのがきっかけで生まれた合同誌です。\par
「わたしも即売会に行ってみたいんですよね。若いうちに一度『人は創作物にここま\par
で本気になれるんだ』っていうのを体験しておきたい、というか。こういうのってそ\par
の後の創作活動の原動力になると思うんですよ」\par
そんな軽口を「なら、出展側で参加してみませんか？」と後押しし、同人誌を作るの\par
も即売会に参加するのもはじめてのわたしにたくさんの助言と支援をくださった方々\par
と、こんなにも頼りないわたしを信じて原稿を託してくださった寄稿者の方々のおか\par
げで、わたしはこうしてみなさんに本誌をお届けできています。本誌がみなさんの記\par
憶に残る一冊になることを願うと同時に、わたしにとっても一生忘れられない一冊に\par
なるだろうと強く感じています。\par
最後に、本誌の刊行に関わったすべての寄稿者の方、外部協力者の方、ならびに『学\par
園アイドルマスター』ひいては『アイドルマスター』シリーズに携わるすべての方、\par
特に篠澤広さんにこの場を借りてお礼申し上げます。\par
2025年8月　サークル『Silent//Azure』代表\par
橘いおね(@TachibanaIone)\par
主宰あとがき\par
62\par

% --- PDF page 64 ---
奥付\par
完全理解！篠澤広の機械学習$\cdot$数理最適化ゼミ合同～1限目～\par
寄　稿 橘いおね\par
Najiko\par
イカ爆弾 / 山船\par
うにかじき\par
betty\par
tokio\par
連絡先 ione@tachibanai.one\par
x.com/TachibanaIone\par
校　正 橘いおね betty\par
編　集 橘いおね\par
発行者 Silent//Azure\par
2025年8月18日 v\textbf{[缺字]}.\textbf{[缺字]}.\textbf{[缺字]}発行\par

% --- PDF page 65 ---
\noindent\textit{[第 65 页没有可抽取文字层。]}\par
\endgroup
